% !TEX program = pdflatex
%
% Top-level entry. Section bodies live under paper/sections/, named so that
% directory sort order matches the document order. Long-only material is
% guarded by \iflong ... \fi; the default build is the long version.
% A short companion entry point (paper_short.tex) overrides \longfalse.

\documentclass{article}

\usepackage[preprint]{neurips_2024}
\usepackage{amsmath,amssymb}
\usepackage{booktabs}
\usepackage{graphicx}
\usepackage{xcolor}
\usepackage{hyperref}
\hypersetup{
  colorlinks=true,
  linkcolor={blue!50!black},   % internal refs (sections, equations, figures)
  citecolor={blue!50!black},   % bibliography citations
  urlcolor={blue!60!black},    % external URLs
}
\usepackage[capitalize, nameinlink]{cleveref}
\usepackage{listings}
\usepackage{caption}

% Long-version flag. paper_short.tex sets \longfalse before % !TEX program = pdflatex
%
% Top-level entry. Section bodies live under paper/sections/, named so that
% directory sort order matches the document order. Long-only material is
% guarded by \iflong ... \fi; the default build is the long version.
% A short companion entry point (paper_short.tex) overrides \longfalse.

\documentclass{article}

\usepackage[preprint]{neurips_2024}
\usepackage{amsmath,amssymb}
\usepackage{booktabs}
\usepackage{graphicx}
\usepackage{xcolor}
\usepackage{hyperref}
\hypersetup{
  colorlinks=true,
  linkcolor={blue!50!black},   % internal refs (sections, equations, figures)
  citecolor={blue!50!black},   % bibliography citations
  urlcolor={blue!60!black},    % external URLs
}
\usepackage[capitalize, nameinlink]{cleveref}
\usepackage{listings}
\usepackage{caption}

% Long-version flag. paper_short.tex sets \longfalse before % !TEX program = pdflatex
%
% Top-level entry. Section bodies live under paper/sections/, named so that
% directory sort order matches the document order. Long-only material is
% guarded by \iflong ... \fi; the default build is the long version.
% A short companion entry point (paper_short.tex) overrides \longfalse.

\documentclass{article}

\usepackage[preprint]{neurips_2024}
\usepackage{amsmath,amssymb}
\usepackage{booktabs}
\usepackage{graphicx}
\usepackage{xcolor}
\usepackage{hyperref}
\hypersetup{
  colorlinks=true,
  linkcolor={blue!50!black},   % internal refs (sections, equations, figures)
  citecolor={blue!50!black},   % bibliography citations
  urlcolor={blue!60!black},    % external URLs
}
\usepackage[capitalize, nameinlink]{cleveref}
\usepackage{listings}
\usepackage{caption}

% Long-version flag. paper_short.tex sets \longfalse before % !TEX program = pdflatex
%
% Top-level entry. Section bodies live under paper/sections/, named so that
% directory sort order matches the document order. Long-only material is
% guarded by \iflong ... \fi; the default build is the long version.
% A short companion entry point (paper_short.tex) overrides \longfalse.

\documentclass{article}

\usepackage[preprint]{neurips_2024}
\usepackage{amsmath,amssymb}
\usepackage{booktabs}
\usepackage{graphicx}
\usepackage{xcolor}
\usepackage{hyperref}
\hypersetup{
  colorlinks=true,
  linkcolor={blue!50!black},   % internal refs (sections, equations, figures)
  citecolor={blue!50!black},   % bibliography citations
  urlcolor={blue!60!black},    % external URLs
}
\usepackage[capitalize, nameinlink]{cleveref}
\usepackage{listings}
\usepackage{caption}

% Long-version flag. paper_short.tex sets \longfalse before \input{paper}
% so a single source can produce both targets.
\providecommand{\iflong}{}
\ifx\iflong\empty
  \newif\iflong
  \longtrue
\fi

% Stand-in markers. Anything tagged \todo{...} is a placeholder for content
% that has been sketched but not written, or where the author still needs to
% confirm a number, citation, or figure. Render in the margin so the gaps are
% visible at a glance during review.
\newcommand{\todo}[1]{\textcolor{red}{\textbf{[TODO: #1]}}}
\newcommand{\sketch}[1]{\textcolor{gray}{\textit{(sketch: #1)}}}
\newcommand{\fig}[1]{\textcolor{blue}{\textbf{[FIG: #1]}}}

\title{Differentiable Extracellular Stimulation
of Compartmental Neuron Models}
\author{%
  Marco Christiani \\
  Carnegie Mellon University \\
  \texttt{marcochr@andrew.cmu.edu}
}

\graphicspath{{./}{figures/}}

\begin{document}
\maketitle

\begin{abstract}
We extend Jaxley, a JAX-native compartmental neuron simulator, with
extracellular stimulation. The extension constructs the discrete activating
function from Jaxley's pre-normalized axial diffusion operator and injects it
as an equivalent compartmental current through the public stimulus API; no
fork or solver patch is required. Every step is JAX-traceable, so
\texttt{jit}, \texttt{vmap}, \texttt{grad}, and \texttt{jax.sharding} compose
with the pipeline without modification. We verify the operator against
Jaxley's internal voltage vector field to $\sim\!10^{-13}$ relative error in
float64 and recover $O(\Delta x^2)$ convergence to the analytical activating
function on a uniform straight cable. We then exercise the pipeline on
Hodgkin--Huxley cables and on the full-morphology Blue Brain Project
L2/3 pyramidal cell, with the parvalbumin basket cell included only at
the single-compartment channel-kinetics level since its full
morphology does not run stably in our setup. We validate against
NEURON in a three-step cross-check that reaches a median per-site
voltage RMSE of $0.446\,$mV on the full-morphology Pyr cell under
extracellular stimulation, sweep stimulation parameters across pulse
width, polarity, distance, and frequency on a four-chip TPU pod-slice,
and report a $\sim\!34\times$ wall-clock speedup over serial
single-core NEURON on the same Pyr cell. Gradients
of voltage observables flow through the full simulator with respect to
electrode position, waveform parameters, and tissue conductivity,
placing direct waveform optimisation within reach of the same compute
budget.
\end{abstract}

\input{sections/01_intro}
\input{sections/02_background}
\input{sections/03_method}
\input{sections/04_cells}
\input{sections/05_verification}
\input{sections/06_results}
\input{sections/07_discussion}
\input{sections/08_limitations}
\input{sections/09_conclusion}

\bibliographystyle{plain}
\bibliography{ref}

\appendix
\input{sections/10_appendix_branchpoint}
\input{sections/11_appendix_equivalent_current}
\input{sections/12_appendix_bbp_impl}
\input{sections/13_appendix_ih_gradient}
\input{sections/14_appendix_reproducibility}

\end{document}

% so a single source can produce both targets.
\providecommand{\iflong}{}
\ifx\iflong\empty
  \newif\iflong
  \longtrue
\fi

% Stand-in markers. Anything tagged \todo{...} is a placeholder for content
% that has been sketched but not written, or where the author still needs to
% confirm a number, citation, or figure. Render in the margin so the gaps are
% visible at a glance during review.
\newcommand{\todo}[1]{\textcolor{red}{\textbf{[TODO: #1]}}}
\newcommand{\sketch}[1]{\textcolor{gray}{\textit{(sketch: #1)}}}
\newcommand{\fig}[1]{\textcolor{blue}{\textbf{[FIG: #1]}}}

\title{Differentiable Extracellular Stimulation
of Compartmental Neuron Models}
\author{%
  Marco Christiani \\
  Carnegie Mellon University \\
  \texttt{marcochr@andrew.cmu.edu}
}

\graphicspath{{./}{figures/}}

\begin{document}
\maketitle

\begin{abstract}
We extend Jaxley, a JAX-native compartmental neuron simulator, with
extracellular stimulation. The extension constructs the discrete activating
function from Jaxley's pre-normalized axial diffusion operator and injects it
as an equivalent compartmental current through the public stimulus API; no
fork or solver patch is required. Every step is JAX-traceable, so
\texttt{jit}, \texttt{vmap}, \texttt{grad}, and \texttt{jax.sharding} compose
with the pipeline without modification. We verify the operator against
Jaxley's internal voltage vector field to $\sim\!10^{-13}$ relative error in
float64 and recover $O(\Delta x^2)$ convergence to the analytical activating
function on a uniform straight cable. We then exercise the pipeline on
Hodgkin--Huxley cables and on the full-morphology Blue Brain Project
L2/3 pyramidal cell, with the parvalbumin basket cell included only at
the single-compartment channel-kinetics level since its full
morphology does not run stably in our setup. We validate against
NEURON in a three-step cross-check that reaches a median per-site
voltage RMSE of $0.446\,$mV on the full-morphology Pyr cell under
extracellular stimulation, sweep stimulation parameters across pulse
width, polarity, distance, and frequency on a four-chip TPU pod-slice,
and report a $\sim\!34\times$ wall-clock speedup over serial
single-core NEURON on the same Pyr cell. Gradients
of voltage observables flow through the full simulator with respect to
electrode position, waveform parameters, and tissue conductivity,
placing direct waveform optimisation within reach of the same compute
budget.
\end{abstract}


% =============================================================
\section{Introduction}\label{sec:intro}
% =============================================================

Clinical neurostimulation devices, including deep brain stimulators and
spinal cord stimulators, deliver pulsed extracellular currents through
electrodes placed near the target tissue. Therapeutic outcome depends on a
small set of waveform parameters --- pulse width, shape, amplitude,
frequency, polarity --- whose effect on neural activation is nonlinear and
sensitive to local geometry. In present clinical practice these parameters
are tuned by hand, guided by symptom observation across patient visits.
Computational design has been explored as an
alternative, but the simulators in widespread use are slow when the goal is
to sweep the design space.

The standard pipeline pairs a volume-conductor field model with the NEURON
simulator~\cite{Hines2001NeuronAT} to integrate the resulting
transmembrane response. The default BBP NEURON workflow runs serially
on a single CPU core, so existing studies that search for energy-efficient or selective
waveforms have either used genetic
algorithms~\cite{grill2015model,wongsarnpigoon2010energy} or trained
surrogate neural networks on NEURON
outputs~\cite{melissano2024efficient}. Genetic search returns no gradient
and requires many forward simulations per design. Surrogates are fast at
inference but introduce approximation error and have to be retrained
whenever the underlying biophysical model changes.

Jaxley~\cite{deistler2025jaxley} is a recently published JAX-native
compartmental simulator. It supports automatic differentiation, JIT
compilation, \texttt{vmap}, and the modern \texttt{jax.sharding} API for
multi-device execution. It does not currently expose extracellular
stimulation. We add it without forking the simulator: the extension
constructs the discrete activating function from Jaxley's own axial
diffusion operator and injects the result as a compartmental current
through the public stimulus API.

\paragraph{Contributions.}
\begin{enumerate}
  \item A differentiable extracellular-stimulation pipeline through
  Jaxley's public stimulus API, with no fork of the upstream simulator.
  \item A derivation of the discrete activating function
  $G\boldsymbol{\phi}$ from Jaxley's pre-normalised voltage diffusion
  operator $G$.
  \item Numerical verification: agreement with Jaxley's internal voltage
  vector field to $\sim\!10^{-13}$ relative error in float64,
  $O(\Delta x^2)$ convergence to the analytical activating function, and
  end-to-end gradient flow through \texttt{jx.integrate}.
  \item Three-step cross-validation against NEURON --- single-compartment
  kinetics, full-morphology Pyr intracellular, and full-morphology Pyr
  under extracellular stimulation --- with characterised residuals at
  each step.
  \item TPU-scale empirical results on the BBP L2/3 pyramidal cell:
  a strength--duration sweep across distance and polarity, a
  charge-balanced biphasic grid of $315$ valid configurations, and a
  throughput--scaling curve at $\sim\!34\times$ wall-clock speedup over
  serial NEURON at $B = 1000$.
\end{enumerate}


% =============================================================
\section{Background and related work}\label{sec:background}
% =============================================================

\subsection{Extracellular potential from point-source electrodes}

At neural stimulation frequencies, displacement currents in tissue are
negligible compared with conduction currents, and Maxwell's equations
reduce to a quasi-static formulation in which the extracellular potential
satisfies Laplace's equation away from sources. Modelling each electrode
contact as a point source in an infinite, homogeneous, isotropic
conductor of conductivity $\sigma$ gives the closed-form monopole
$\phi(\mathbf{x}) = I / (4\pi\sigma\,d)$, with $d$ the source-to-field
distance~\cite{10.1093/acprof:oso/9780195058239.001.0001,10.1093/acprof:oso/9780195050387.001.0001}.
Multiple electrodes follow by superposition. We use this model throughout; the homogeneity and isotropy
assumptions are limitations we return to in \cref{sec:limitations}.

\subsection{The activating function}

For a uniform straight cable in an extracellular field $\phi(x,t)$,
Rattay~\cite{rattay1986analysis} showed that the driving term for
transmembrane voltage is proportional to
$\partial^2 \phi / \partial x^2$ along the cable; McNeal's earlier
compartmental treatment of myelinated axons under monopolar
stimulation~\cite{Mcneal1976AnalysisOA} contains the same idea in
discrete form. On an arbitrary compartmental tree the discrete
analogue is $[G\boldsymbol{\phi}]_j$, where $G$ is the axial diffusion
operator that the cable solver already maintains; the continuous second
derivative is the special case of a uniform unbranched cable. We exploit
this in \cref{sec:method}: rather than recomputing
$\partial^2 \phi / \partial x^2$ along each branch, we apply Jaxley's
$G$ directly to the sampled extracellular potential.

\subsection{Differentiable neuron simulators}

Jaxley~\cite{deistler2025jaxley} is the only published JAX-native
compartmental simulator we are aware of that supports HH-style channel
kinetics, branched morphologies, and reverse-mode autodiff through the
time integrator. Diffrax~\cite{DBLP:journals/corr/abs-2202-02435} provides
differentiable ODE solvers in JAX, but it does not model compartmental
neurons. Brian2~\cite{Stimberg2019} is
Python-native and considerably faster than NEURON, but it is not built
for automatic differentiation. We reuse Jaxley's solver through its public
stimulus API rather than forking the simulator: the extracellular extension
is a JAX module that composes with the rest of the framework without
modifying its internals.


% =============================================================
\section{Method}\label{sec:method}
% =============================================================

We adopt the notation of \cref{tab:notation}. The pipeline is:
\begingroup
\setlength{\abovedisplayskip}{4pt}
\setlength{\belowdisplayskip}{4pt}
\begin{equation}
  I_e(t) \xrightarrow{\text{field}} \boldsymbol{\Phi}
  \xrightarrow{G\,\cdot} \mathbf{f}
  \xrightarrow{\text{encode}} \mathbf{I}_{\text{ecs}}
  \xrightarrow{\text{integrate}} \mathbf{v}(t).
  \label{eq:pipeline}
\end{equation}
\endgroup

\begin{table}[t]
\centering
\small
\caption{Notation. Lengths in $\mu$m, currents in $\mu$A or nA, voltages
in mV, time in ms, conductivity in S/m.}
\label{tab:notation}
\begin{tabular}{@{}lll@{}}
\toprule
Symbol & Meaning & Units \\ \midrule
$N$ & Number of compartments & --- \\
$T$ & Number of timesteps & --- \\
$\mathbf{x}_j,\,\mathbf{x}_e$ & Compartment center, electrode position & $\mu$m \\
$d_{je}$ & $\lVert \mathbf{x}_j - \mathbf{x}_e \rVert$ & $\mu$m \\
$I_e(t)$ & Electrode current (cathodic negative) & $\mu$A \\
$\sigma$ & Extracellular conductivity & S/m \\
$\phi_j(t),\,\boldsymbol{\Phi}$ & Extracellular potential, $\boldsymbol{\Phi}_{jt}=\phi_j(t)$ & mV \\
$v_j(t),\,\mathbf{v}$ & Transmembrane voltage & mV \\
$G$ & Axial voltage diffusion operator & $1/$ms \\
$c_j,\,A_j$ & Specific capacitance, surface area & $\mu$F/cm$^2$,\,$\mu$m$^2$ \\
$f_j(t) = [G\boldsymbol{\phi}(t)]_j$ & Activating function at compartment $j$ & mV/ms \\
\bottomrule
\end{tabular}
\end{table}

\paragraph{Extracellular potential from point sources.}
For an electrode $e$ delivering $I_e(t)$ in a homogeneous, isotropic,
infinite volume conductor of conductivity $\sigma$, superposition of
monopoles gives
\[
  \phi_j(t) \;=\; \sum_e \frac{I_e(t)}{4\pi\sigma\,d_{je}}.
\]
Defining the transfer factor $h_{je} = 10^3 / (4\pi\sigma\,d_{je})$
(in $\mathrm{mV}/\mu\mathrm{A}$, with the $10^3$ absorbing the unit
conversion from $\mu$A and $\mu$m to mV) factorises space and time, so
that for a single electrode
$\boldsymbol{\Phi} = \mathbf{h}\,\mathbf{I}^\top \in \mathbb{R}^{N\times T}$.

\paragraph{The discrete activating function.}
Jaxley solves the spatially discretised cable equation
\[
  c_j A_j \frac{dv_j}{dt}
  = \sum_{k\sim j} g_{jk}(v_k - v_j) - A_j I_{\text{ion}}(v_j,\mathbf{s}_j) + i_{\text{ext},j},
\]
where $g_{jk}$ is the raw axial conductance (units S) between
compartments $j$ and $k$, $c_j$ is specific membrane capacitance, and
$A_j$ is compartment surface area. Internally, Jaxley folds $c_j A_j$
and unit conversions into a single voltage diffusion operator $G$ with
entries $\tilde g_{jk} = g_{jk} / (c_j A_j)$ (in $1/$ms; note the
asymmetric row-$j$ normalisation, which is why $G$ is symmetric only
when $c_j A_j$ is uniform across compartments), so that
$\dot{\mathbf{v}} = G\mathbf{v} + (\text{membrane terms})$. The sign
pattern of $G$ is the weighted graph Laplacian:
$G_{jk} = +\tilde g_{jk}$ for adjacent $k\neq j$,
$G_{jj} = -\sum_{k\neq j} \tilde g_{jk}$, and zero otherwise. Rows sum
to zero (current conservation).

An extracellular field $\phi_j(t)$ shifts the axial driving force from
$v_k - v_j$ to $(v_k+\phi_k)-(v_j+\phi_j)$; by linearity of $G$, this
is equivalent to adding $G\boldsymbol{\phi}$ as a forcing term, the
discrete analogue of Rattay's continuous activating function
$\rho_a^{-1}\partial^2\phi/\partial x^2$~\cite{rattay1986analysis}.
On a uniform straight cable with constant radius and compartment
length $\Delta x$, $[G\boldsymbol{\phi}]_j$ at interior compartments
reduces to a constant multiple of the standard
$[1,-2,1]/\Delta x^2$ stencil for $\partial^2\phi/\partial x^2$, with
$O(\Delta x^2)$ truncation error (verified in
\cref{sec:results-verify}).
On branched morphologies, we delegate the elimination of the
branchpoint pseudo-nodes that Jaxley introduces at cable junctions to
its own solver routine, keeping the operator in sync with upstream
solver changes (\cref{app:branchpoint-elimination}).

\paragraph{Equivalent injected current.}
The activating function $G\boldsymbol{\phi}$ has units mV/ms (a voltage
rate), but Jaxley's stimulus API accepts current in nA. We encode the
forcing as a per-compartment current $\mathbf{I}_{\text{ecs}}$ chosen
so that, after Jaxley's internal current-to-voltage-rate conversion,
the integrator sees exactly $G\boldsymbol{\phi}$. The unit-conversion
algebra is given in \cref{app:equivalent-current}; the membrane
capacitance and area factors that appear in the encoding cancel
exactly in the round trip and are an artifact of working through the
public API rather than modifying the solver.

\paragraph{Implementation.}
The pipeline of \cref{eq:pipeline} is realised as four small JAX-native
modules wired together by an experiment object.
\texttt{field.point\_source\_potential} computes $\boldsymbol{\Phi}$ from
electrode positions, currents, and $\sigma$, with a $1\,\mu$m floor on
$d_{je}$ to regularise the $1/d$ singularity when an electrode lands inside
a compartment.
\texttt{discretization.build\_voltage\_operator\_G} returns $G$ by
delegating to Jaxley's \texttt{\_compute\_transition\_matrix} (the same
routine the simulator uses for its backward-Euler step) and stripping
branchpoint pseudo-nodes.
\texttt{equivalent\_current.phi\_e\_to\_ecs\_nA} applies the encoding
above.
An \texttt{ECSExperiment} class precomputes $G$, $\mathbf{c}$, $\mathbf{A}$,
and the compartment centres in its constructor, and exposes
\texttt{simulate\_waveform} for a single integration and
\texttt{simulate\_and\_extract} for one integration plus response-feature
extraction.

All operations are JAX-traceable: \texttt{jit}, \texttt{vmap},
\texttt{grad}, and \texttt{jax.sharding.NamedSharding} compose with
the pipeline without modification, so a one-dimensional device mesh is
enough to data-parallelise the batch axis across a TPU pod-slice. Two
implementation choices matter at scale and unlock the full-fidelity
throughput regime of \cref{sec:res-throughput}: $G$ is stored as a
sparse \texttt{BCOO} matrix to avoid the $(B,\,N,\,N)$ XLA broadcast
that dense $G$ triggers under \texttt{vmap}, and
\texttt{simulate\_waveform} forwards a \texttt{checkpoint\_lengths}
factorisation of $T$ into \texttt{jx.integrate} so the backward-Euler
scan rematerialises in $\sim\!\sqrt{T}$ chunks rather than holding the
full $(T,\,B,\,N)$ voltage history in HBM.


% =============================================================
\section{Cell models}\label{sec:cells}
% =============================================================

\subsection{Hodgkin--Huxley straight cable}\label{sec:cells-hh}

The verification studies of \cref{sec:results-verify} and the pilot
strength--duration sweep of \cref{sec:res-sd} use a uniform
Hodgkin--Huxley cable. Defaults reproduce the configuration set by
\texttt{experiment.make\_hh\_cable\_experiment}: $50$ compartments
along a $1.25\,$mm cable, fibre radius $10\,\mu$m, axial resistivity
$\rho_a = 100\,\Omega\cdot$cm, with the standard Jaxley HH channel set on
every compartment. A single electrode sits $50\,\mu$m above the first
compartment in a medium of conductivity $\sigma = 0.3\,$S/m, the typical
value for cortical grey matter.

\subsection{BBP L2/3 cells (cADpyr229, cNAC187)}\label{sec:cells-bbp}

For morphologically detailed simulations we adapt two L2/3 cell models
from the Blue Brain Project~\cite{markram2015reconstruction}: the
regular-spiking pyramidal cell
\texttt{L23\_PC\_cADpyr229} and the parvalbumin-expressing basket cell
\texttt{L23\_LBC\_cNAC187}. Channel kinetics are reused from
\texttt{jaxley\_mech.channels.l5pc}, which provides Jaxley-compatible
implementations of the BBP L5 channel set; conductances and passive
parameters are read from the BBP biophysics specification and applied per
section group (soma, axon, apical, basal).

Two implementation details affect reproducibility but do not bear on
the paper's results: the SWC conversion of the BBP soma contour, and
the calcium-channel insertion order in Jaxley. Both are documented in
\cref{app:bbp-implementation}.

A third detail is needed for NEURON parity. The \texttt{cADpyr229}
apical $I_h$ conductance follows a distance-dependent gradient
specified in BBP's biophysics file with a path-distance argument; we
apply it under a Euclidean-distance approximation to the soma centre,
which closes a $0.6\,$mV resting-potential offset that a uniform
somatic value otherwise leaves (closed-form gradient and
implementation in \cref{app:ih-gradient}). Long SWC branches are
subdivided at load time (\texttt{max\_branch\_len}$\,=\,100\,\mu$m) so
the compartment count tracks what NEURON's \texttt{geom\_nseg} routine
produces for the same morphology; we found this necessary for
spike-timing parity (\cref{sec:discussion}).


% =============================================================
\section{Numerical verification}\label{sec:results-verify}
% =============================================================

Three numerical checks anchor the rest of the paper: that $G$ matches what
Jaxley's solver computes internally, that $G\boldsymbol{\phi}$ recovers the
continuous activating function as the discretisation refines, and that
gradients propagate through the full simulator. All three are exercised in
the test suite (\texttt{test\_ecs\_operator\_consistency} and
\texttt{test\_ecs\_grad}).

\paragraph{Operator round-trip.}
We compare $G\mathbf{v}$ against Jaxley's internal axial term on a
50-compartment cable in float64, isolating the axial term by zeroing
the membrane contributions (\texttt{voltage\_terms}\,$=\,0$ and
\texttt{constant\_terms}\,$=\,0$). The two outputs agree to
$\sim\!10^{-13}$ per-compartment relative error
(\cref{fig:verification}, left); once $G\mathbf{v}$ matches Jaxley's
own axial term, $G\boldsymbol{\phi}$ is the correct extracellular
forcing on the same tree.

\paragraph{$O(\Delta x^2)$ convergence.}
For a uniform straight cable with a point-source electrode the activating
function has the closed form $\rho_a^{-1}\,\partial^2\phi/\partial x^2$.
We compare $G\boldsymbol{\phi}$ against this analytical reference on the
central $10\,\%$ of the cable (boundary effects dominate near the ends)
across compartment counts $\{50,\,100,\,200,\,400\}$. The relative error
decreases by a factor of $\sim 4$ per doubling
($3.95\times,\,3.99\times,\,4.00\times$), consistent with the
$O(\Delta x^2)$ rate of the underlying $[1,-2,1]/\Delta x^2$ stencil
(\cref{fig:verification}, right).

\begin{figure}[t]
  \centering
  \includegraphics[width=\linewidth]{figures/verification.png}
  \caption{Numerical verification of the discrete activating function.
  Left: histogram of per-compartment relative error of $G\mathbf{v}$ vs
  Jaxley's internal \texttt{\_voltage\_vectorfield} on a $50$-compartment
  HH cable in float64; the median is $\sim\!3\!\times\!10^{-16}$ and no
  compartment exceeds $\sim\!10^{-13}$. Right: relative error of
  $G\boldsymbol{\phi}$ versus the closed-form
  $\rho_a^{-1}\partial^2\phi/\partial x^2$ on the central $10\,\%$ of the
  cable, plotted against $\Delta x$ on log--log axes; the dashed line is
  a reference slope-2, and the box reports the per-doubling error
  ratios.}
  \label{fig:verification}
\end{figure}

\paragraph{Gradient flow.}
The pipeline is JAX-traceable end to end, so \texttt{jax.grad} of any
voltage observable propagates through the field model, the operator,
the encoding, and Jaxley's integrator. Unit tests confirm finite,
non-zero gradients with respect to electrode position, waveform, and
tissue conductivity. At full BBP scale, the position gradient through
a $700$-compartment Pyr cell under a $-50\,\mu$A, $1\,$ms cathodic
pulse from electrode $(0, 100, 0)\,\mu$m is
\[
  \frac{\partial v_{\mathrm{soma,\,peak}}}{\partial \mathbf{x}_e}
  = (\,0.330,\,-0.371,\,-0.243\,)\;\mathrm{mV}/\mu\mathrm{m},
\]
returned in $37.6\,$s on a v5e single chip. The full receipt is part
of the artifact set described in the reproducibility appendix.


% =============================================================
\section{Empirical results}\label{sec:results-empirical}
% =============================================================

This section answers four questions that a stimulation designer or
methodologist would put to the simulator. Does it agree with NEURON on
a realistic morphology before we trust any extracellular result? At a
given electrode distance and pulse width, what amplitude is needed to
evoke a spike? Where in the asymmetric biphasic waveform space does
the cell respond, given the same charge per pulse? And does the
infrastructure scale to design-space sweeps that NEURON cannot run? We
close with a gradient receipt that motivates the next step --- direct
optimisation of the waveform itself. All runs reported below are part
of the curated artifact set described in the reproducibility appendix;
wall-clock figures are on a single v5e chip unless stated otherwise.

\subsection{Validation against NEURON}\label{sec:res-neuron}

Before any stimulation result from this pipeline is trustworthy, the
simulator has to agree with NEURON on the same cell. We validate in
three steps, each adding one new mechanism on top of the previous so
that residuals can be attributed.

\paragraph{Setup and confounders.}
The NEURON reference uses the same BBP cADpyr229 morphology that
Jaxley loads, with three confounders equalised before measurement.
(i) Axon myelination is enabled by default in the BBP NEURON template;
we disable it for the comparison since our Jaxley implementation does
not model myelin. (ii) The \texttt{jaxley\_mech.channels.l5pc} ports
apply a fixed $Q_{10}$ temperature correction; we run NEURON at
$34\,^\circ$C to match. (iii) NEURON's \texttt{geom\_nseg} produces $1131$
segments on this morphology, while Jaxley with \texttt{ncomp}$\,=\,2$
and \texttt{max\_branch\_len}$\,=\,100$ produces $700$ compartments;
we report the mismatch rather than forcing them to match.

\paragraph{Step 1: channel kinetics.} We reduce the morphology to a
$20\,\mu$m$\times 20\,\mu$m cylinder (\texttt{parity\_single\_soma}),
so that the comparison exercises only channel kinetics and ion
reversals. With the full BBP Pyr and PV channel sets installed, spike
counts agree to $\pm 1$ over $\sim\!50$ spikes per cell (Pyr: $18$
NEURON / $16$ Jaxley; PV: $50$/$51$), and AP peak heights agree to
$\pm 0.04\,$mV across both cells. The match is consistent with a
source-level diff of the relevant \texttt{jaxley\_mech} channels
against the BBP \texttt{.mod} files.

\paragraph{Step 2: full-morphology Pyr, intracellular.} We instantiate
the full cADpyr229 morphology in both simulators and inject a
$0.5\,$nA somatic step current with extracellular coupling held out
(\texttt{parity\_bbp\_pyr}). Both simulators emit $3$ spikes
(\cref{fig:parity-intracellular}); first-spike latency differs by
$-50\,\mu$s ($2$ integration steps at $\Delta t = 0.025\,$ms), and AP
peak by $+0.49\,$mV. Subthreshold
waveform RMSE is $0.275\,$mV with Pearson $r = 0.9995$;
suprathreshold full-trace RMSE is $6.032\,$mV with $r = 0.9253$, where
the latter is dominated by spike-edge timing offsets. The
$-0.38\,$mV resting-potential residual is explained by our
Euclidean-distance approximation to BBP's path-distance apical $I_h$
gradient (\cref{sec:cells-bbp}); a path-distance traversal would close
it further.

\paragraph{Step 3: full-morphology Pyr, extracellular.} The headline
claim of this paper is that ECS works on a realistic branched cell.
We test it by comparing voltage traces with NEURON when both
simulators see the same point-source field on the same Pyr
morphology. After aligning Jaxley's compartment centres and NEURON's
segment readouts to the same physical sites
(\texttt{parity\_ecs.bbp\_pyr}), per-site voltage RMSE is
$0.387,\,0.553,\,0.505,\,0.322\,$mV at soma, apical, basal and axon
respectively (median $0.446\,$mV); the site-local $\phi$ matches
between the two simulators (\cref{fig:parity-ecs}). The ECS residuals should be read against
the Step-2 baseline, since the underlying solver and discretisation
differences carry over. Combined with the analytical
$O(\Delta x^2)$ convergence of \cref{sec:results-verify} on a uniform
straight cable, this anchors the activating-function pipeline at both
geometric extremes: closed-form correctness on a straight cable and
external parity on a branched real cell. The practical consequence is
that ECS-on-Pyr in this pipeline reproduces NEURON behaviour to
sub-millivolt residuals on a realistic morphology --- enough fidelity
to support design-space sweeps that NEURON cannot run in the same time
budget (\cref{sec:res-throughput}).

\begin{figure}[t]
  \centering
  \includegraphics[width=\linewidth]{figures/bbp_intracellular_parity.png}
  \caption{Step 2 (full-morphology Pyr, intracellular) cross-check.
  NEURON (black) and Jaxley (dashed) soma voltage traces under $0.1\,$nA
  subthreshold and $0.5\,$nA suprathreshold step currents on the same
  cADpyr229 morphology. Subthreshold RMSE $0.275\,$mV
  ($r = 0.9995$); suprathreshold RMSE $6.032\,$mV
  ($r = 0.9253$), dominated by spike-edge timing offsets
  attributable to the $700$-vs-$1131$ compartmentalisation difference.}
  \label{fig:parity-intracellular}
\end{figure}

\begin{figure}[t]
  \centering
  \includegraphics[width=\linewidth]{figures/bbp_ecs_parity.png}
  \caption{Step 3 (full-morphology Pyr, extracellular) cross-check.
  NEURON (black) and Jaxley (dashed) voltage traces at the soma and
  apical readout sites after aligning the two simulators' compartment
  centres to within $0\,\mu$m on the same Pyr morphology, plus per-site
  RMSE bars and the site-local $\phi$ that drives the ECS coupling.
  Median per-site RMSE $0.446\,$mV; the residuals carry over from the
  Step-2 baseline rather than reflecting a new ECS-specific
  discrepancy.}
  \label{fig:parity-ecs}
\end{figure}

\subsection{Strength--duration on the HH cable}\label{sec:res-sd}

At a given electrode distance and pulse width, what amplitude does it
take to evoke a spike? The strength--duration curve is the standard
first answer to that question, and underlies almost every clinical and
engineering decision downstream of it~\cite{4122088}. We define the threshold $A^*$ for a fixed
configuration $(d, \text{polarity}, t_p)$ as the smallest electrode
amplitude that elicits a single soma spike in the simulation:
\[
  A^*(d,\,\text{polarity},\,t_p)
  \;=\; \inf\bigl\{A \ge 0\,:\,F(w_{A,t_p};\,d) \text{ spikes}\bigr\}.
\]
Each $A^*$ is found by a one-dimensional binary search in $A$, with the
spike-detection criterion threshold-crossing of the soma voltage (see
\texttt{response.detect\_spike}). The full sweep is one binary search per
$(d,\,\text{polarity},\,t_p)$ triple --- four electrode distances
$\{25,\,50,\,100,\,200\}\,\mu$m, three waveform polarities (monophasic
cathodic, monophasic anodic, charge-balanced biphasic), and nine pulse
widths from $0.05$ to $2\,$ms --- for $108$ thresholds in total.

The search uses a uniform $[0,\,5000]\,\mu$A bracket with $14$
bisection iterations, giving a resolution of $\sim\!0.3\,\mu$A; spike
detection stays stable across this range for the HH cable. The whole
sweep runs in float-32. Cathodic and biphasic thresholds are well-resolved at
every $(d,\,t_p)$ pair and follow the textbook strength--duration
shape, declining toward a rheobase as $t_p$ grows
(\cref{fig:strength-duration}). Monophasic anodic elicits spikes only
at short pulse widths and short distances; it saturates at the bracket
ceiling for longer pulses. This polarity asymmetry is the discrete
analogue of the activating-function sign argument: cathodic stimulation
sets up a depolarising forcing in the central segments where the cable
fires first, whereas anodic stimulation depolarises the cable ends,
which only triggers under brief, geometrically-favourable
conditions~\cite{rattay1986analysis,Mcneal1976AnalysisOA}. For a
designer choosing parameters, the takeaway is that
ms-scale cathodic pulses minimise charge per spike at every distance
in the sweep, charge-balanced biphasic costs roughly two-fold more,
and pure anodic stimulation is only viable in the short-pulse corner.

\begin{figure}[t]
  \centering
  \includegraphics[width=\linewidth]{figures/strength_duration.png}
  \caption{Strength--duration curves on the $50$-compartment HH cable.
  One panel per electrode distance; one curve per waveform polarity.
  Threshold amplitude $|A^*|$ falls toward a rheobase as pulse width
  grows, and grows with electrode distance as the $1/d$ field decay
  implies. The dotted line is the binary-search ceiling ($5000\,\mu$A);
  pulse-width points missing from a curve indicate no spike below that
  ceiling, which characterises the anodic regime at long widths.}
  \label{fig:strength-duration}
\end{figure}

\subsection{Charge-balanced biphasic grid}\label{sec:res-grid}

Asymmetric biphasic pulses are the clinical lever for trading
activation selectivity against charge efficiency: at the same total
delivered charge per pulse, where on the $(T_p, T_n)$ plane does the
cell actually respond? A symmetric charge-balanced biphasic pulse has
two equal-and-opposite phases of equal duration; a clinical waveform
almost always has unequal phase durations (longer-low cathodic +
short-high anodic, or vice versa)~\cite{wongsarnpigoon2010energy}. To characterise
this region of waveform space, we hold the electrode at $d=100\,\mu$m
(the cleanest distance from the strength--duration sweep) and grid four
parameters: anodic phase duration $T_p$, cathodic phase duration $T_n$,
anodic amplitude $A_p$ (with cathodic amplitude $A_n$ pinned by
$A_p T_p = A_n T_n$ to keep total injected charge zero), and pulse-train
frequency $f$. We record spike count, latency, and peak voltage per
configuration rather than searching for a threshold, so each grid point
is one forward simulation.

The grid has $T_p,\,T_n \in \{50,\,100,\,200,\,500\}\,\mu$s
(four values each), $A_p \in \{10,\,20,\,30,\,40,\,80\}\,\mu$A, and
pulse-train frequencies $f \in \{100,\,200,\,500,\,1000\}\,$Hz, for
$315$ valid charge-balanced forward simulations; the five invalid
combinations have $T_p + T_n$ equal to the $1000\,$Hz period. The full grid
runs in $16\,$min on the v5litepod-4 pod-slice via
\texttt{vmap}-over-configuration with hierarchical scan-checkpointing
(\cref{sec:res-throughput}). Spike crossings concentrate along the
asymmetric $T_p = 500\,\mu$s edge at low-to-mid amplitudes
($A_p \le 30\,\mu$A), and at $A_p = 80\,\mu$A for $T_p = 200\,\mu$s
(\cref{fig:biphasic-grid}); the symmetric diagonal $T_p = T_n$ is far
less effective at the same charge per pulse. The
implication for waveform design is direct: asymmetry is not a marginal
parameter --- on this cell, at this distance, swapping a symmetric
biphasic pulse for an asymmetric one with the same charge can be the
difference between sub-threshold and reliable spiking.

\begin{figure}[t]
  \centering
  \includegraphics[width=\linewidth]{figures/biphasic_grid.png}
  \caption{Charge-balanced biphasic grid on the BBP L2/3 Pyr cell at
  $d = 100\,\mu$m. Each panel is a $(T_p, T_n)$ heatmap of the maximum
  soma voltage over the four pulse-train frequencies, with one panel
  per anodic amplitude $A_p$. The dotted diagonal $T_p = T_n$ marks the
  symmetric-biphasic line; cells where $\mathit{vmax}$ crosses
  $0\,$mV correspond to a soma spike at some frequency in the sweep.
  White cells indicate configurations for which all simulated
  frequencies returned non-finite voltage traces and therefore no
  interpretable $\mathit{vmax}$; these occur mainly in the
  long-$T_p$ high-amplitude regime. Five period-filling configurations
  ($T_p = T_n = 500\,\mu$s at $1000\,$Hz) are absent from the
  forward-simulation set.}
  \label{fig:biphasic-grid}
\end{figure}

\subsection{Throughput scaling on TPU}\label{sec:res-throughput}

The point of the simulator extension is to support design-space
sweeps that were previously infeasible. We measure how it scales
across two regimes that probe different bottlenecks. A
\emph{cheap-cells} sweep isolates per-cell compile and kernel-dispatch
overhead at small problem size: \texttt{ncomp}$=2$ on the BBP Pyr
($\sim\!700$ compartments per cell after branchpoint expansion), $B$
up to $1024$ on the v5litepod-4 pod-slice. A \emph{full-fidelity}
sweep is the actual scientific workload: $50$-compartment branches on
the same Pyr cell ($\approx\!17\,500$ compartments per cell), $B$ up
to $4096$ on the same pod-slice. Each $B$ value is a single
\texttt{jax.jit(jax.vmap(\ldots))} dispatch; we report the cached
(JIT-warm) timing.

The cheap-cells curve scales linearly to $B=1024$ at $5.1\,$ms per cell
($\sim\!195$ cells/s); the full-fidelity curve scales linearly through
$B=1024$, peaks at $7.28$ cells per second at $B=2048$, and rolls over
at $B=4096$ as the four chips approach arithmetic-intensity saturation
(\cref{fig:throughput-scaling}). The full-fidelity regime depends on
the sparse-$G$ and \texttt{checkpoint\_lengths} choices in the Method
(\cref{sec:method}); without either, the $(T,\,B,\,N)$ voltage
history overflows the v5e's $16\,$GB chip memory at $B \ge 256$.

\begin{figure}[t]
  \centering
  \includegraphics[width=\linewidth]{figures/throughput_v5p4_scaling.png}
  \caption{Jaxley batch throughput on the v5litepod-4 pod-slice for two
  workload regimes: cheap cells ($700$ compartments per cell, blue) and
  full-fidelity BBP Pyr ($\sim\!17\,500$ compartments per cell, red).
  Throughput scales linearly through $B=1024$ in both regimes; the
  full-fidelity curve peaks at $7.28$ cells per second at $B=2048$ and
  saturates at $B=4096$. JIT-warm timings.}
  \label{fig:throughput-scaling}
\end{figure}

\paragraph{Comparison with NEURON.} The same $700$-compartment Pyr cell
on a $T_\text{sim} = 200\,$ms protocol runs at $0.286$ simulations per
second under serial single-core NEURON (the default BBP setup), against
$9.67$ simulations per second for Jaxley at $B = 1000$ on a single TPU
once the JIT cache is warm --- a $\sim\!34\times$ wall-clock speedup
(\cref{fig:throughput}).
NEURON has parallel-CPU variants such as
CoreNEURON~\cite{Kumbhar2019CoreNEURONA} and
experimental GPU support that we did not benchmark; the relevant
comparison for our purposes is between a CPU-serial reference and an
accelerator-batched formulation. AI accelerators are SIMT/SIMD
throughput devices, and a batched-\texttt{vmap} formulation is the
workload pattern they are designed for; a CPU-bound serial framework
cannot expose the same batch parallelism without a similar
reformulation. The practical consequence: the parameter sweeps in
\cref{sec:res-sd,sec:res-grid} finish in tens of minutes, putting
sweeps over the full waveform-and-geometry design space within reach
of a single TPU job.

\begin{figure}[t]
  \centering
  \includegraphics[width=\linewidth]{figures/bbp_throughput.png}
  \caption{Wall-clock comparison on the same $700$-compartment Pyr cell
  and $T_\text{sim} = 200\,$ms protocol: serial single-core NEURON
  (blue) versus Jaxley on a single TPU chip (red). At $B=1000$, NEURON
  takes $58.3\,$min (measured), Jaxley takes $1.7\,$min ($103.4\,$ms
  per simulation, JIT-warm), giving a $\sim\!34\times$ wall-clock
  speedup. The dashed blue line is the linear extrapolation of the
  per-NEURON-sim cost.}
  \label{fig:throughput}
\end{figure}

\subsection{Gradient-based design signal}\label{sec:res-grad}

For the differentiability we built into the pipeline to matter for
stimulation design, the gradients have to point in the right
direction --- and their direction has to mean something physical.
Establishing that gradients exist (\cref{sec:results-verify}) is a
prerequisite for, but not equivalent to, showing that they carry usable
design signal. The position gradient reported in
\cref{sec:results-verify} admits a direct physical reading. With
the electrode at $(0,\,100,\,0)\,\mu$m, the negative $y$-component of
$\partial v_\mathrm{soma,\,peak} / \partial \mathbf{x}_e$ means moving
the electrode farther from the cell along its dorsal axis reduces the
soma response, consistent with the $1/d$ decay in the field model; the
horizontal components measure the local field anisotropy induced by the
asymmetric placement of dendrites relative to the electrode. Free-form
optimisation of $w(t)$ against a spike-elicitation target under
charge-balance and amplitude constraints is a natural next step and is
left as future work (\cref{sec:limitations}). What this section
establishes is that the same pipeline that runs the empirical sweeps
of \cref{sec:res-sd,sec:res-grid} also returns design-relevant
gradients in a single backward pass --- placing direct waveform
optimisation, rather than gradient-free search, within reach of the
same compute budget.

\begin{figure}[t]
  \centering
  \includegraphics[width=0.78\linewidth]{figures/gradient_arrow.png}
  \caption{Position-gradient direction on the BBP L2/3 Pyr cell. Grey
  points are the $700$-compartment morphology projected onto the
  $y$-$z$ plane; the soma is marked with a black dot, the electrode
  ($\mathbf{x}_e = (0,\,100,\,0)\,\mu$m) with a black cross. The
  orange arrow shows the direction of
  $\partial v_\mathrm{soma,\,peak} / \partial \mathbf{x}_e$, i.e.\ the
  direction in which moving the electrode increases the soma peak.}
  \label{fig:gradient-arrow}
\end{figure}

% =============================================================
\section{Discussion}\label{sec:discussion}
% =============================================================

Genetic search over
waveforms~\cite{grill2015model,wongsarnpigoon2010energy} returns no
gradient direction for a candidate that is close to but short of an
activation target; surrogate networks trained on NEURON
outputs~\cite{hussain2024efficient} introduce an approximation error
that must be re-incurred whenever the underlying biophysical model
changes. Because the pipeline is JAX-traceable end-to-end, gradients
$\partial \ell / \partial w(t)$ of any differentiable scalar objective
$\ell$ with respect to the input waveform are returned in a single
backward pass and avoid both costs. As a
sharper receipt than the position gradient of \cref{sec:res-grad}, an
Adam loop on $\mathrm{MSE}(\hat{v}, v_\mathrm{target})$ on the full
BBP Pyr morphology recovers \texttt{gNaTs2T} from a $-15\,\%$
perturbation to within $+0.13\,\%$ in nine steps
(\texttt{fit\_demo}); the point of this fit is end-to-end gradient
flow on a realistic cell, not optimiser performance. A selectivity
Pareto sweep against a non-target cell, or a charge-minimising waveform
under a spike-elicitation constraint, follows the same template (a
differentiable spike-elicitation surrogate plus charge-balance and
amplitude constraints).

% =============================================================
\section{Limitations and future work}\label{sec:limitations}
% =============================================================

\paragraph{Limitations.}
\begin{itemize}
  \item \emph{Volume conductor.} The field model is a homogeneous,
  isotropic, infinite point-source monopole. Tissue anisotropy,
  finite electrode geometry (cuff, paddle, contact size), and the
  bone--CSF--grey-matter layering used in DBS planning are not
  represented.
  \item \emph{Quasi-static assumption.} Tissue capacitance is ignored;
  this is standard at neural-stimulation frequencies but breaks down at
  very short pulse edges.
  \item \emph{Float-32 caveat.} The production sweeps run in float-32
  after a $0.77\,\%$ threshold-difference check against float-64 on one
  HH cable configuration; the check should be repeated per cell type
  before reusing float-32 in a regime we have not validated.
  \item \emph{PV cell stability.} The full-morphology BBP L2/3 LBC
  basket cell is unstable in Jaxley's default integrator (NaN at
  $\sim\!30\,$ms even unstimulated), traced to BBP's
  $g_\mathrm{pas} = 1\!\times\!10^{-6}\,$S/cm$^2$ on PV apical and
  basal --- $30\times$ lower than the Pyr value and tolerated by NEURON
  but not by the Jaxley defaults. Single-compartment PV parity passes
  (\cref{sec:res-neuron}); all full-morphology runtime claims in this
  paper are therefore Pyr-only.
  \item \emph{Pyr Ih path-distance traversal.} The
  $-0.38\,$mV resting-potential offset against NEURON
  (\cref{sec:res-neuron}) is the residual of approximating BBP's
  path-distance apical $I_h$ gradient with Euclidean distance.
  Replacing the approximation with a parent-pointer breadth-first
  traversal would close the residual but we did not implement it.
  \item \emph{Solver floor.} \texttt{checkpoint\_lengths} bounds memory
  but not arithmetic. The backward-Euler scan dominates the per-step
  cost on a v5e chip and is what saturates the throughput curve at
  $B=2048$ in \cref{sec:res-throughput}.
\end{itemize}

\paragraph{Future work.}
Three lines of follow-up sit immediately on top of what is here. First,
restoring full PV cell coverage by stabilising the LBC integration ---
either with an explicit smaller-timestep regime, a leak-conductance
floor, or a stiff-aware solver --- and applying the BBP path-distance
$I_h$ gradient to PV apical and basal as in Pyr. Second, gradient-based
optimisation of the free-form waveform $w(t)$ against a
spike-elicitation surrogate, under charge-balance and amplitude
constraints, evaluated against the strength--duration baseline of
\cref{sec:res-sd}. Third, selectivity: two cell models with different
geometries placed in the same extracellular field, optimising $w(t)$
to drive a soma spike in the target cell while keeping the non-target
subthreshold. The Pareto frontier between activation and sparing is
exactly what the differentiable pipeline is built to trace; a
penalty-based formulation
$\mathcal{L}(w) = \lVert w \rVert_1 \Delta t
+ \lambda\,\tilde{F}(w;\theta_2)$ subject to
$\tilde{F}(w;\theta_1) \ge \tau$ keeps the problem smooth and lets a
$\lambda$-sweep trace the frontier directly.


% =============================================================
\section{Conclusion}\label{sec:conclusion}
% =============================================================

We have added extracellular stimulation to Jaxley by recovering the
discrete activating function from the simulator's own axial diffusion
operator and injecting it as an equivalent compartmental current through
the public stimulus API. The pipeline composes with \texttt{jit},
\texttt{vmap}, \texttt{grad}, and \texttt{jax.sharding} without a fork
of the simulator, agrees with Jaxley's internal voltage vector field to
$\sim\!10^{-13}$ relative error, and runs the BBP L2/3 pyramidal cell at
full $50$-compartment-per-branch fidelity on a four-chip v5e pod-slice.
What this enables, and what we leave to follow-up work, is gradient-based
search over the waveform itself --- the design problem that motivated
the extension.



\bibliographystyle{plain}
\bibliography{ref}

\appendix

% =============================================================
\section{Branchpoint elimination}\label{app:branchpoint-elimination}
% =============================================================

Jaxley's internal representation introduces branchpoint pseudo-nodes
at cable junctions to enforce Kirchhoff's current law. Partitioning
$G$ into real ($R$) and branchpoint ($B$) blocks gives
\[
  \begin{pmatrix} G_{RR} & G_{RB} \\ G_{BR} & G_{BB} \end{pmatrix}
  \begin{pmatrix} \mathbf{v}_R \\ \mathbf{v}_B \end{pmatrix}
  =
  \begin{pmatrix} [G\mathbf{v}]_R \\ [G\mathbf{v}]_B \end{pmatrix}.
\]
Branchpoint nodes carry no membrane, so $c_b A_b = 0$ at $b$ and the
cable equation at $b$ reduces to the algebraic constraint
$[G\mathbf{v}]_B = 0$, giving $\mathbf{v}_B = -G_{BB}^{-1}G_{BR}\mathbf{v}_R$
and the Schur complement
\[
  \tilde{G} \;=\; G_{RR} - G_{RB}\,G_{BB}^{-1}\,G_{BR}.
\]
The same operator applies to the extracellular forcing:
$\tilde{G}\boldsymbol{\phi}_R$ gives the activating function at the
real compartments, where $\boldsymbol{\phi}_R$ is the field sampled at
real compartment centres only. The reduced operator acts only on real
compartments, so no $\boldsymbol{\phi}_B$ value enters the forcing. Our implementation delegates this assembly to a
private Jaxley helper rather than reconstructing it from scratch, so
the operator stays in sync with upstream solver changes.


% =============================================================
\section{Equivalent-current encoding}\label{app:equivalent-current}
% =============================================================

The activating function $G\boldsymbol{\phi}$ has units mV/ms (a
voltage rate), but Jaxley's stimulus API accepts current in nA.
Internally Jaxley converts
$i_{\text{ext}}\,[\mu\text{A/cm}^2] = i_{\text{nA}}\,[\text{nA}] /
A\,[\mu\text{m}^2]\cdot 10^5$ and adds $i_{\text{ext}}/c$ to $dv/dt$.
To inject $[G\boldsymbol{\phi}]_j$ on the other side of this
conversion, we encode
\[
  i_{\text{ecs},j}\,[\text{nA}] \;=\;
    c_j \cdot [G\boldsymbol{\phi}(t)]_j \cdot \frac{A_j}{10^5}.
\]
The $c_j$ and $A_j$ factors cancel in the round trip; they are an
encoding artifact of using the public API rather than modifying the
solver. In matrix form,
\[
  \mathbf{I}_{\text{ecs}} \;=\;
    \left(\mathbf{c} \odot \frac{\mathbf{A}}{10^5}\right)
    \odot (G\boldsymbol{\Phi})
    \;\in\; \mathbb{R}^{N\times T}.
\]


% =============================================================
\section{BBP cell-model implementation notes}\label{app:bbp-implementation}
% =============================================================

Two implementation issues affect reproducibility of the BBP cell models
in Jaxley but are unrelated to the paper's results.

\paragraph{SWC soma conversion.} The BBP source morphologies represent
the soma as a two-dimensional contour, which the SWC format cannot
encode. The Neurolucida-to-SWC conversion tool emits soma points with
$r = 0$, which Jaxley turns into degenerate zero-area compartments
that produce NaN voltages at the first integration step. We patch the
converted SWC files at load time so the soma points carry a
non-degenerate radius.

\paragraph{Calcium-channel insertion order.} Calcium signalling is
sensitive to channel insertion order. The calcium Nernst reversal
must be established before any calcium-carrying channel, which in
turn must precede the calcium pump and the calcium-activated
potassium channel.
\texttt{cell\_factory.\_insert\_calcium\_chain} enforces
\texttt{CaNernstReversal} $\to$ \texttt{CaHVA/CaLVA} $\to$
\texttt{CaPump} $\to$ \texttt{SKE2}; out-of-order insertion produces
uninitialised state references at integration time.


% =============================================================
\section{Apical \texorpdfstring{$I_h$}{Ih} distance-dependent
gradient}\label{app:ih-gradient}
% =============================================================

The \texttt{cADpyr229} apical $I_h$ conductance follows BBP's
distance-dependent gradient
\[
  \bar{g}_{I_h}(d)
    = \bigl(-0.8696 + 2.087\,e^{0.0031\,d}\bigr)
      \cdot 8\!\times\!10^{-5}\;\mathrm{S/cm}^2,
\]
where $d$ is path distance from the soma centre in $\mu$m. We
approximate $d$ with the Euclidean distance to the soma centre, which
is close for roughly-vertical apical trees. Before the fix the
resting-potential offset against NEURON was $-0.62\,$mV; the
Euclidean approximation closes it to $-0.38\,$mV, and a path-distance
traversal would close it further (\cref{sec:limitations}).


% =============================================================
\section{Reproducibility}\label{app:reproducibility}
% =============================================================

The build pipeline binds the entire compute environment, not just the
Python packages. A \texttt{requirements.txt} inside a Dockerfile pins
Python packages and a base OS image but leaves the C and
C\nolinkurl{++} compilers, the CUDA toolchain, the Python interpreter
build, and the cloud topology unspecified --- each of which can drift
between machines and change numerical results. The stack we use
closes those gaps.

\paragraph{Python dependencies (\texttt{uv}).} A Rust-implemented
package manager that resolves direct and transitive Python
dependencies and records every resulting wheel and its hash in a
lockfile (\texttt{uv.lock}). Two machines sharing the lockfile
install the same versions of every Python package, regardless of
install order.

\paragraph{System environment (Nix).} A Nix flake (\texttt{flake.nix})
specifies the Python interpreter (3.11), CUDA runtime, JAX, the C
compiler that built them, and command-line tooling. Store paths are
determined by the hash of the full transitive build graph (pinned via
\texttt{flake.lock}), so the same flake on the same target system
produces the same dependency closure.

\paragraph{Cloud infrastructure (OpenTofu).} The TPU pod-slice, Cloud
SQL Postgres database, and GCS artifact bucket are declared as
OpenTofu configuration in \texttt{infra/tofu/}; \texttt{tofu apply}
brings the topology up, \texttt{tofu destroy} tears it down.

\paragraph{Reproducing the figures.} A reader entering the Nix shell
gets a deterministic Python and CUDA environment; the project
\texttt{README} carries the exact commands. Reproducing the figures
from pinned data does not require any TPU-scale workload.

\paragraph{Artifact set.} The curated data backing every figure ---
parity \texttt{.npz} files, the strength--duration and biphasic-grid
Zarr stores, the throughput sweeps, and the gradient receipt ---
totals a few megabytes and is published as a separate artifact
distribution rather than committed to the source repository, available
at \url{https://github.com/Marco-Christiani/jaxley-extracellular/releases/tag/paper-v0.1}.


\end{document}

% so a single source can produce both targets.
\providecommand{\iflong}{}
\ifx\iflong\empty
  \newif\iflong
  \longtrue
\fi

% Stand-in markers. Anything tagged \todo{...} is a placeholder for content
% that has been sketched but not written, or where the author still needs to
% confirm a number, citation, or figure. Render in the margin so the gaps are
% visible at a glance during review.
\newcommand{\todo}[1]{\textcolor{red}{\textbf{[TODO: #1]}}}
\newcommand{\sketch}[1]{\textcolor{gray}{\textit{(sketch: #1)}}}
\newcommand{\fig}[1]{\textcolor{blue}{\textbf{[FIG: #1]}}}

\title{Differentiable Extracellular Stimulation
of Compartmental Neuron Models}
\author{%
  Marco Christiani \\
  Carnegie Mellon University \\
  \texttt{marcochr@andrew.cmu.edu}
}

\graphicspath{{./}{figures/}}

\begin{document}
\maketitle

\begin{abstract}
We extend Jaxley, a JAX-native compartmental neuron simulator, with
extracellular stimulation. The extension constructs the discrete activating
function from Jaxley's pre-normalized axial diffusion operator and injects it
as an equivalent compartmental current through the public stimulus API; no
fork or solver patch is required. Every step is JAX-traceable, so
\texttt{jit}, \texttt{vmap}, \texttt{grad}, and \texttt{jax.sharding} compose
with the pipeline without modification. We verify the operator against
Jaxley's internal voltage vector field to $\sim\!10^{-13}$ relative error in
float64 and recover $O(\Delta x^2)$ convergence to the analytical activating
function on a uniform straight cable. We then exercise the pipeline on
Hodgkin--Huxley cables and on the full-morphology Blue Brain Project
L2/3 pyramidal cell, with the parvalbumin basket cell included only at
the single-compartment channel-kinetics level since its full
morphology does not run stably in our setup. We validate against
NEURON in a three-step cross-check that reaches a median per-site
voltage RMSE of $0.446\,$mV on the full-morphology Pyr cell under
extracellular stimulation, sweep stimulation parameters across pulse
width, polarity, distance, and frequency on a four-chip TPU pod-slice,
and report a $\sim\!34\times$ wall-clock speedup over serial
single-core NEURON on the same Pyr cell. Gradients
of voltage observables flow through the full simulator with respect to
electrode position, waveform parameters, and tissue conductivity,
placing direct waveform optimisation within reach of the same compute
budget.
\end{abstract}


% =============================================================
\section{Introduction}\label{sec:intro}
% =============================================================

Clinical neurostimulation devices, including deep brain stimulators and
spinal cord stimulators, deliver pulsed extracellular currents through
electrodes placed near the target tissue. Therapeutic outcome depends on a
small set of waveform parameters --- pulse width, shape, amplitude,
frequency, polarity --- whose effect on neural activation is nonlinear and
sensitive to local geometry. In present clinical practice these parameters
are tuned by hand, guided by symptom observation across patient visits.
Computational design has been explored as an
alternative, but the simulators in widespread use are slow when the goal is
to sweep the design space.

The standard pipeline pairs a volume-conductor field model with the NEURON
simulator~\cite{Hines2001NeuronAT} to integrate the resulting
transmembrane response. The default BBP NEURON workflow runs serially
on a single CPU core, so existing studies that search for energy-efficient or selective
waveforms have either used genetic
algorithms~\cite{grill2015model,wongsarnpigoon2010energy} or trained
surrogate neural networks on NEURON
outputs~\cite{melissano2024efficient}. Genetic search returns no gradient
and requires many forward simulations per design. Surrogates are fast at
inference but introduce approximation error and have to be retrained
whenever the underlying biophysical model changes.

Jaxley~\cite{deistler2025jaxley} is a recently published JAX-native
compartmental simulator. It supports automatic differentiation, JIT
compilation, \texttt{vmap}, and the modern \texttt{jax.sharding} API for
multi-device execution. It does not currently expose extracellular
stimulation. We add it without forking the simulator: the extension
constructs the discrete activating function from Jaxley's own axial
diffusion operator and injects the result as a compartmental current
through the public stimulus API.

\paragraph{Contributions.}
\begin{enumerate}
  \item A differentiable extracellular-stimulation pipeline through
  Jaxley's public stimulus API, with no fork of the upstream simulator.
  \item A derivation of the discrete activating function
  $G\boldsymbol{\phi}$ from Jaxley's pre-normalised voltage diffusion
  operator $G$.
  \item Numerical verification: agreement with Jaxley's internal voltage
  vector field to $\sim\!10^{-13}$ relative error in float64,
  $O(\Delta x^2)$ convergence to the analytical activating function, and
  end-to-end gradient flow through \texttt{jx.integrate}.
  \item Three-step cross-validation against NEURON --- single-compartment
  kinetics, full-morphology Pyr intracellular, and full-morphology Pyr
  under extracellular stimulation --- with characterised residuals at
  each step.
  \item TPU-scale empirical results on the BBP L2/3 pyramidal cell:
  a strength--duration sweep across distance and polarity, a
  charge-balanced biphasic grid of $315$ valid configurations, and a
  throughput--scaling curve at $\sim\!34\times$ wall-clock speedup over
  serial NEURON at $B = 1000$.
\end{enumerate}


% =============================================================
\section{Background and related work}\label{sec:background}
% =============================================================

\subsection{Extracellular potential from point-source electrodes}

At neural stimulation frequencies, displacement currents in tissue are
negligible compared with conduction currents, and Maxwell's equations
reduce to a quasi-static formulation in which the extracellular potential
satisfies Laplace's equation away from sources. Modelling each electrode
contact as a point source in an infinite, homogeneous, isotropic
conductor of conductivity $\sigma$ gives the closed-form monopole
$\phi(\mathbf{x}) = I / (4\pi\sigma\,d)$, with $d$ the source-to-field
distance~\cite{10.1093/acprof:oso/9780195058239.001.0001,10.1093/acprof:oso/9780195050387.001.0001}.
Multiple electrodes follow by superposition. We use this model throughout; the homogeneity and isotropy
assumptions are limitations we return to in \cref{sec:limitations}.

\subsection{The activating function}

For a uniform straight cable in an extracellular field $\phi(x,t)$,
Rattay~\cite{rattay1986analysis} showed that the driving term for
transmembrane voltage is proportional to
$\partial^2 \phi / \partial x^2$ along the cable; McNeal's earlier
compartmental treatment of myelinated axons under monopolar
stimulation~\cite{Mcneal1976AnalysisOA} contains the same idea in
discrete form. On an arbitrary compartmental tree the discrete
analogue is $[G\boldsymbol{\phi}]_j$, where $G$ is the axial diffusion
operator that the cable solver already maintains; the continuous second
derivative is the special case of a uniform unbranched cable. We exploit
this in \cref{sec:method}: rather than recomputing
$\partial^2 \phi / \partial x^2$ along each branch, we apply Jaxley's
$G$ directly to the sampled extracellular potential.

\subsection{Differentiable neuron simulators}

Jaxley~\cite{deistler2025jaxley} is the only published JAX-native
compartmental simulator we are aware of that supports HH-style channel
kinetics, branched morphologies, and reverse-mode autodiff through the
time integrator. Diffrax~\cite{DBLP:journals/corr/abs-2202-02435} provides
differentiable ODE solvers in JAX, but it does not model compartmental
neurons. Brian2~\cite{Stimberg2019} is
Python-native and considerably faster than NEURON, but it is not built
for automatic differentiation. We reuse Jaxley's solver through its public
stimulus API rather than forking the simulator: the extracellular extension
is a JAX module that composes with the rest of the framework without
modifying its internals.


% =============================================================
\section{Method}\label{sec:method}
% =============================================================

We adopt the notation of \cref{tab:notation}. The pipeline is:
\begingroup
\setlength{\abovedisplayskip}{4pt}
\setlength{\belowdisplayskip}{4pt}
\begin{equation}
  I_e(t) \xrightarrow{\text{field}} \boldsymbol{\Phi}
  \xrightarrow{G\,\cdot} \mathbf{f}
  \xrightarrow{\text{encode}} \mathbf{I}_{\text{ecs}}
  \xrightarrow{\text{integrate}} \mathbf{v}(t).
  \label{eq:pipeline}
\end{equation}
\endgroup

\begin{table}[t]
\centering
\small
\caption{Notation. Lengths in $\mu$m, currents in $\mu$A or nA, voltages
in mV, time in ms, conductivity in S/m.}
\label{tab:notation}
\begin{tabular}{@{}lll@{}}
\toprule
Symbol & Meaning & Units \\ \midrule
$N$ & Number of compartments & --- \\
$T$ & Number of timesteps & --- \\
$\mathbf{x}_j,\,\mathbf{x}_e$ & Compartment center, electrode position & $\mu$m \\
$d_{je}$ & $\lVert \mathbf{x}_j - \mathbf{x}_e \rVert$ & $\mu$m \\
$I_e(t)$ & Electrode current (cathodic negative) & $\mu$A \\
$\sigma$ & Extracellular conductivity & S/m \\
$\phi_j(t),\,\boldsymbol{\Phi}$ & Extracellular potential, $\boldsymbol{\Phi}_{jt}=\phi_j(t)$ & mV \\
$v_j(t),\,\mathbf{v}$ & Transmembrane voltage & mV \\
$G$ & Axial voltage diffusion operator & $1/$ms \\
$c_j,\,A_j$ & Specific capacitance, surface area & $\mu$F/cm$^2$,\,$\mu$m$^2$ \\
$f_j(t) = [G\boldsymbol{\phi}(t)]_j$ & Activating function at compartment $j$ & mV/ms \\
\bottomrule
\end{tabular}
\end{table}

\paragraph{Extracellular potential from point sources.}
For an electrode $e$ delivering $I_e(t)$ in a homogeneous, isotropic,
infinite volume conductor of conductivity $\sigma$, superposition of
monopoles gives
\[
  \phi_j(t) \;=\; \sum_e \frac{I_e(t)}{4\pi\sigma\,d_{je}}.
\]
Defining the transfer factor $h_{je} = 10^3 / (4\pi\sigma\,d_{je})$
(in $\mathrm{mV}/\mu\mathrm{A}$, with the $10^3$ absorbing the unit
conversion from $\mu$A and $\mu$m to mV) factorises space and time, so
that for a single electrode
$\boldsymbol{\Phi} = \mathbf{h}\,\mathbf{I}^\top \in \mathbb{R}^{N\times T}$.

\paragraph{The discrete activating function.}
Jaxley solves the spatially discretised cable equation
\[
  c_j A_j \frac{dv_j}{dt}
  = \sum_{k\sim j} g_{jk}(v_k - v_j) - A_j I_{\text{ion}}(v_j,\mathbf{s}_j) + i_{\text{ext},j},
\]
where $g_{jk}$ is the raw axial conductance (units S) between
compartments $j$ and $k$, $c_j$ is specific membrane capacitance, and
$A_j$ is compartment surface area. Internally, Jaxley folds $c_j A_j$
and unit conversions into a single voltage diffusion operator $G$ with
entries $\tilde g_{jk} = g_{jk} / (c_j A_j)$ (in $1/$ms; note the
asymmetric row-$j$ normalisation, which is why $G$ is symmetric only
when $c_j A_j$ is uniform across compartments), so that
$\dot{\mathbf{v}} = G\mathbf{v} + (\text{membrane terms})$. The sign
pattern of $G$ is the weighted graph Laplacian:
$G_{jk} = +\tilde g_{jk}$ for adjacent $k\neq j$,
$G_{jj} = -\sum_{k\neq j} \tilde g_{jk}$, and zero otherwise. Rows sum
to zero (current conservation).

An extracellular field $\phi_j(t)$ shifts the axial driving force from
$v_k - v_j$ to $(v_k+\phi_k)-(v_j+\phi_j)$; by linearity of $G$, this
is equivalent to adding $G\boldsymbol{\phi}$ as a forcing term, the
discrete analogue of Rattay's continuous activating function
$\rho_a^{-1}\partial^2\phi/\partial x^2$~\cite{rattay1986analysis}.
On a uniform straight cable with constant radius and compartment
length $\Delta x$, $[G\boldsymbol{\phi}]_j$ at interior compartments
reduces to a constant multiple of the standard
$[1,-2,1]/\Delta x^2$ stencil for $\partial^2\phi/\partial x^2$, with
$O(\Delta x^2)$ truncation error (verified in
\cref{sec:results-verify}).
On branched morphologies, we delegate the elimination of the
branchpoint pseudo-nodes that Jaxley introduces at cable junctions to
its own solver routine, keeping the operator in sync with upstream
solver changes (\cref{app:branchpoint-elimination}).

\paragraph{Equivalent injected current.}
The activating function $G\boldsymbol{\phi}$ has units mV/ms (a voltage
rate), but Jaxley's stimulus API accepts current in nA. We encode the
forcing as a per-compartment current $\mathbf{I}_{\text{ecs}}$ chosen
so that, after Jaxley's internal current-to-voltage-rate conversion,
the integrator sees exactly $G\boldsymbol{\phi}$. The unit-conversion
algebra is given in \cref{app:equivalent-current}; the membrane
capacitance and area factors that appear in the encoding cancel
exactly in the round trip and are an artifact of working through the
public API rather than modifying the solver.

\paragraph{Implementation.}
The pipeline of \cref{eq:pipeline} is realised as four small JAX-native
modules wired together by an experiment object.
\texttt{field.point\_source\_potential} computes $\boldsymbol{\Phi}$ from
electrode positions, currents, and $\sigma$, with a $1\,\mu$m floor on
$d_{je}$ to regularise the $1/d$ singularity when an electrode lands inside
a compartment.
\texttt{discretization.build\_voltage\_operator\_G} returns $G$ by
delegating to Jaxley's \texttt{\_compute\_transition\_matrix} (the same
routine the simulator uses for its backward-Euler step) and stripping
branchpoint pseudo-nodes.
\texttt{equivalent\_current.phi\_e\_to\_ecs\_nA} applies the encoding
above.
An \texttt{ECSExperiment} class precomputes $G$, $\mathbf{c}$, $\mathbf{A}$,
and the compartment centres in its constructor, and exposes
\texttt{simulate\_waveform} for a single integration and
\texttt{simulate\_and\_extract} for one integration plus response-feature
extraction.

All operations are JAX-traceable: \texttt{jit}, \texttt{vmap},
\texttt{grad}, and \texttt{jax.sharding.NamedSharding} compose with
the pipeline without modification, so a one-dimensional device mesh is
enough to data-parallelise the batch axis across a TPU pod-slice. Two
implementation choices matter at scale and unlock the full-fidelity
throughput regime of \cref{sec:res-throughput}: $G$ is stored as a
sparse \texttt{BCOO} matrix to avoid the $(B,\,N,\,N)$ XLA broadcast
that dense $G$ triggers under \texttt{vmap}, and
\texttt{simulate\_waveform} forwards a \texttt{checkpoint\_lengths}
factorisation of $T$ into \texttt{jx.integrate} so the backward-Euler
scan rematerialises in $\sim\!\sqrt{T}$ chunks rather than holding the
full $(T,\,B,\,N)$ voltage history in HBM.


% =============================================================
\section{Cell models}\label{sec:cells}
% =============================================================

\subsection{Hodgkin--Huxley straight cable}\label{sec:cells-hh}

The verification studies of \cref{sec:results-verify} and the pilot
strength--duration sweep of \cref{sec:res-sd} use a uniform
Hodgkin--Huxley cable. Defaults reproduce the configuration set by
\texttt{experiment.make\_hh\_cable\_experiment}: $50$ compartments
along a $1.25\,$mm cable, fibre radius $10\,\mu$m, axial resistivity
$\rho_a = 100\,\Omega\cdot$cm, with the standard Jaxley HH channel set on
every compartment. A single electrode sits $50\,\mu$m above the first
compartment in a medium of conductivity $\sigma = 0.3\,$S/m, the typical
value for cortical grey matter.

\subsection{BBP L2/3 cells (cADpyr229, cNAC187)}\label{sec:cells-bbp}

For morphologically detailed simulations we adapt two L2/3 cell models
from the Blue Brain Project~\cite{markram2015reconstruction}: the
regular-spiking pyramidal cell
\texttt{L23\_PC\_cADpyr229} and the parvalbumin-expressing basket cell
\texttt{L23\_LBC\_cNAC187}. Channel kinetics are reused from
\texttt{jaxley\_mech.channels.l5pc}, which provides Jaxley-compatible
implementations of the BBP L5 channel set; conductances and passive
parameters are read from the BBP biophysics specification and applied per
section group (soma, axon, apical, basal).

Two implementation details affect reproducibility but do not bear on
the paper's results: the SWC conversion of the BBP soma contour, and
the calcium-channel insertion order in Jaxley. Both are documented in
\cref{app:bbp-implementation}.

A third detail is needed for NEURON parity. The \texttt{cADpyr229}
apical $I_h$ conductance follows a distance-dependent gradient
specified in BBP's biophysics file with a path-distance argument; we
apply it under a Euclidean-distance approximation to the soma centre,
which closes a $0.6\,$mV resting-potential offset that a uniform
somatic value otherwise leaves (closed-form gradient and
implementation in \cref{app:ih-gradient}). Long SWC branches are
subdivided at load time (\texttt{max\_branch\_len}$\,=\,100\,\mu$m) so
the compartment count tracks what NEURON's \texttt{geom\_nseg} routine
produces for the same morphology; we found this necessary for
spike-timing parity (\cref{sec:discussion}).


% =============================================================
\section{Numerical verification}\label{sec:results-verify}
% =============================================================

Three numerical checks anchor the rest of the paper: that $G$ matches what
Jaxley's solver computes internally, that $G\boldsymbol{\phi}$ recovers the
continuous activating function as the discretisation refines, and that
gradients propagate through the full simulator. All three are exercised in
the test suite (\texttt{test\_ecs\_operator\_consistency} and
\texttt{test\_ecs\_grad}).

\paragraph{Operator round-trip.}
We compare $G\mathbf{v}$ against Jaxley's internal axial term on a
50-compartment cable in float64, isolating the axial term by zeroing
the membrane contributions (\texttt{voltage\_terms}\,$=\,0$ and
\texttt{constant\_terms}\,$=\,0$). The two outputs agree to
$\sim\!10^{-13}$ per-compartment relative error
(\cref{fig:verification}, left); once $G\mathbf{v}$ matches Jaxley's
own axial term, $G\boldsymbol{\phi}$ is the correct extracellular
forcing on the same tree.

\paragraph{$O(\Delta x^2)$ convergence.}
For a uniform straight cable with a point-source electrode the activating
function has the closed form $\rho_a^{-1}\,\partial^2\phi/\partial x^2$.
We compare $G\boldsymbol{\phi}$ against this analytical reference on the
central $10\,\%$ of the cable (boundary effects dominate near the ends)
across compartment counts $\{50,\,100,\,200,\,400\}$. The relative error
decreases by a factor of $\sim 4$ per doubling
($3.95\times,\,3.99\times,\,4.00\times$), consistent with the
$O(\Delta x^2)$ rate of the underlying $[1,-2,1]/\Delta x^2$ stencil
(\cref{fig:verification}, right).

\begin{figure}[t]
  \centering
  \includegraphics[width=\linewidth]{figures/verification.png}
  \caption{Numerical verification of the discrete activating function.
  Left: histogram of per-compartment relative error of $G\mathbf{v}$ vs
  Jaxley's internal \texttt{\_voltage\_vectorfield} on a $50$-compartment
  HH cable in float64; the median is $\sim\!3\!\times\!10^{-16}$ and no
  compartment exceeds $\sim\!10^{-13}$. Right: relative error of
  $G\boldsymbol{\phi}$ versus the closed-form
  $\rho_a^{-1}\partial^2\phi/\partial x^2$ on the central $10\,\%$ of the
  cable, plotted against $\Delta x$ on log--log axes; the dashed line is
  a reference slope-2, and the box reports the per-doubling error
  ratios.}
  \label{fig:verification}
\end{figure}

\paragraph{Gradient flow.}
The pipeline is JAX-traceable end to end, so \texttt{jax.grad} of any
voltage observable propagates through the field model, the operator,
the encoding, and Jaxley's integrator. Unit tests confirm finite,
non-zero gradients with respect to electrode position, waveform, and
tissue conductivity. At full BBP scale, the position gradient through
a $700$-compartment Pyr cell under a $-50\,\mu$A, $1\,$ms cathodic
pulse from electrode $(0, 100, 0)\,\mu$m is
\[
  \frac{\partial v_{\mathrm{soma,\,peak}}}{\partial \mathbf{x}_e}
  = (\,0.330,\,-0.371,\,-0.243\,)\;\mathrm{mV}/\mu\mathrm{m},
\]
returned in $37.6\,$s on a v5e single chip. The full receipt is part
of the artifact set described in the reproducibility appendix.


% =============================================================
\section{Empirical results}\label{sec:results-empirical}
% =============================================================

This section answers four questions that a stimulation designer or
methodologist would put to the simulator. Does it agree with NEURON on
a realistic morphology before we trust any extracellular result? At a
given electrode distance and pulse width, what amplitude is needed to
evoke a spike? Where in the asymmetric biphasic waveform space does
the cell respond, given the same charge per pulse? And does the
infrastructure scale to design-space sweeps that NEURON cannot run? We
close with a gradient receipt that motivates the next step --- direct
optimisation of the waveform itself. All runs reported below are part
of the curated artifact set described in the reproducibility appendix;
wall-clock figures are on a single v5e chip unless stated otherwise.

\subsection{Validation against NEURON}\label{sec:res-neuron}

Before any stimulation result from this pipeline is trustworthy, the
simulator has to agree with NEURON on the same cell. We validate in
three steps, each adding one new mechanism on top of the previous so
that residuals can be attributed.

\paragraph{Setup and confounders.}
The NEURON reference uses the same BBP cADpyr229 morphology that
Jaxley loads, with three confounders equalised before measurement.
(i) Axon myelination is enabled by default in the BBP NEURON template;
we disable it for the comparison since our Jaxley implementation does
not model myelin. (ii) The \texttt{jaxley\_mech.channels.l5pc} ports
apply a fixed $Q_{10}$ temperature correction; we run NEURON at
$34\,^\circ$C to match. (iii) NEURON's \texttt{geom\_nseg} produces $1131$
segments on this morphology, while Jaxley with \texttt{ncomp}$\,=\,2$
and \texttt{max\_branch\_len}$\,=\,100$ produces $700$ compartments;
we report the mismatch rather than forcing them to match.

\paragraph{Step 1: channel kinetics.} We reduce the morphology to a
$20\,\mu$m$\times 20\,\mu$m cylinder (\texttt{parity\_single\_soma}),
so that the comparison exercises only channel kinetics and ion
reversals. With the full BBP Pyr and PV channel sets installed, spike
counts agree to $\pm 1$ over $\sim\!50$ spikes per cell (Pyr: $18$
NEURON / $16$ Jaxley; PV: $50$/$51$), and AP peak heights agree to
$\pm 0.04\,$mV across both cells. The match is consistent with a
source-level diff of the relevant \texttt{jaxley\_mech} channels
against the BBP \texttt{.mod} files.

\paragraph{Step 2: full-morphology Pyr, intracellular.} We instantiate
the full cADpyr229 morphology in both simulators and inject a
$0.5\,$nA somatic step current with extracellular coupling held out
(\texttt{parity\_bbp\_pyr}). Both simulators emit $3$ spikes
(\cref{fig:parity-intracellular}); first-spike latency differs by
$-50\,\mu$s ($2$ integration steps at $\Delta t = 0.025\,$ms), and AP
peak by $+0.49\,$mV. Subthreshold
waveform RMSE is $0.275\,$mV with Pearson $r = 0.9995$;
suprathreshold full-trace RMSE is $6.032\,$mV with $r = 0.9253$, where
the latter is dominated by spike-edge timing offsets. The
$-0.38\,$mV resting-potential residual is explained by our
Euclidean-distance approximation to BBP's path-distance apical $I_h$
gradient (\cref{sec:cells-bbp}); a path-distance traversal would close
it further.

\paragraph{Step 3: full-morphology Pyr, extracellular.} The headline
claim of this paper is that ECS works on a realistic branched cell.
We test it by comparing voltage traces with NEURON when both
simulators see the same point-source field on the same Pyr
morphology. After aligning Jaxley's compartment centres and NEURON's
segment readouts to the same physical sites
(\texttt{parity\_ecs.bbp\_pyr}), per-site voltage RMSE is
$0.387,\,0.553,\,0.505,\,0.322\,$mV at soma, apical, basal and axon
respectively (median $0.446\,$mV); the site-local $\phi$ matches
between the two simulators (\cref{fig:parity-ecs}). The ECS residuals should be read against
the Step-2 baseline, since the underlying solver and discretisation
differences carry over. Combined with the analytical
$O(\Delta x^2)$ convergence of \cref{sec:results-verify} on a uniform
straight cable, this anchors the activating-function pipeline at both
geometric extremes: closed-form correctness on a straight cable and
external parity on a branched real cell. The practical consequence is
that ECS-on-Pyr in this pipeline reproduces NEURON behaviour to
sub-millivolt residuals on a realistic morphology --- enough fidelity
to support design-space sweeps that NEURON cannot run in the same time
budget (\cref{sec:res-throughput}).

\begin{figure}[t]
  \centering
  \includegraphics[width=\linewidth]{figures/bbp_intracellular_parity.png}
  \caption{Step 2 (full-morphology Pyr, intracellular) cross-check.
  NEURON (black) and Jaxley (dashed) soma voltage traces under $0.1\,$nA
  subthreshold and $0.5\,$nA suprathreshold step currents on the same
  cADpyr229 morphology. Subthreshold RMSE $0.275\,$mV
  ($r = 0.9995$); suprathreshold RMSE $6.032\,$mV
  ($r = 0.9253$), dominated by spike-edge timing offsets
  attributable to the $700$-vs-$1131$ compartmentalisation difference.}
  \label{fig:parity-intracellular}
\end{figure}

\begin{figure}[t]
  \centering
  \includegraphics[width=\linewidth]{figures/bbp_ecs_parity.png}
  \caption{Step 3 (full-morphology Pyr, extracellular) cross-check.
  NEURON (black) and Jaxley (dashed) voltage traces at the soma and
  apical readout sites after aligning the two simulators' compartment
  centres to within $0\,\mu$m on the same Pyr morphology, plus per-site
  RMSE bars and the site-local $\phi$ that drives the ECS coupling.
  Median per-site RMSE $0.446\,$mV; the residuals carry over from the
  Step-2 baseline rather than reflecting a new ECS-specific
  discrepancy.}
  \label{fig:parity-ecs}
\end{figure}

\subsection{Strength--duration on the HH cable}\label{sec:res-sd}

At a given electrode distance and pulse width, what amplitude does it
take to evoke a spike? The strength--duration curve is the standard
first answer to that question, and underlies almost every clinical and
engineering decision downstream of it~\cite{4122088}. We define the threshold $A^*$ for a fixed
configuration $(d, \text{polarity}, t_p)$ as the smallest electrode
amplitude that elicits a single soma spike in the simulation:
\[
  A^*(d,\,\text{polarity},\,t_p)
  \;=\; \inf\bigl\{A \ge 0\,:\,F(w_{A,t_p};\,d) \text{ spikes}\bigr\}.
\]
Each $A^*$ is found by a one-dimensional binary search in $A$, with the
spike-detection criterion threshold-crossing of the soma voltage (see
\texttt{response.detect\_spike}). The full sweep is one binary search per
$(d,\,\text{polarity},\,t_p)$ triple --- four electrode distances
$\{25,\,50,\,100,\,200\}\,\mu$m, three waveform polarities (monophasic
cathodic, monophasic anodic, charge-balanced biphasic), and nine pulse
widths from $0.05$ to $2\,$ms --- for $108$ thresholds in total.

The search uses a uniform $[0,\,5000]\,\mu$A bracket with $14$
bisection iterations, giving a resolution of $\sim\!0.3\,\mu$A; spike
detection stays stable across this range for the HH cable. The whole
sweep runs in float-32. Cathodic and biphasic thresholds are well-resolved at
every $(d,\,t_p)$ pair and follow the textbook strength--duration
shape, declining toward a rheobase as $t_p$ grows
(\cref{fig:strength-duration}). Monophasic anodic elicits spikes only
at short pulse widths and short distances; it saturates at the bracket
ceiling for longer pulses. This polarity asymmetry is the discrete
analogue of the activating-function sign argument: cathodic stimulation
sets up a depolarising forcing in the central segments where the cable
fires first, whereas anodic stimulation depolarises the cable ends,
which only triggers under brief, geometrically-favourable
conditions~\cite{rattay1986analysis,Mcneal1976AnalysisOA}. For a
designer choosing parameters, the takeaway is that
ms-scale cathodic pulses minimise charge per spike at every distance
in the sweep, charge-balanced biphasic costs roughly two-fold more,
and pure anodic stimulation is only viable in the short-pulse corner.

\begin{figure}[t]
  \centering
  \includegraphics[width=\linewidth]{figures/strength_duration.png}
  \caption{Strength--duration curves on the $50$-compartment HH cable.
  One panel per electrode distance; one curve per waveform polarity.
  Threshold amplitude $|A^*|$ falls toward a rheobase as pulse width
  grows, and grows with electrode distance as the $1/d$ field decay
  implies. The dotted line is the binary-search ceiling ($5000\,\mu$A);
  pulse-width points missing from a curve indicate no spike below that
  ceiling, which characterises the anodic regime at long widths.}
  \label{fig:strength-duration}
\end{figure}

\subsection{Charge-balanced biphasic grid}\label{sec:res-grid}

Asymmetric biphasic pulses are the clinical lever for trading
activation selectivity against charge efficiency: at the same total
delivered charge per pulse, where on the $(T_p, T_n)$ plane does the
cell actually respond? A symmetric charge-balanced biphasic pulse has
two equal-and-opposite phases of equal duration; a clinical waveform
almost always has unequal phase durations (longer-low cathodic +
short-high anodic, or vice versa)~\cite{wongsarnpigoon2010energy}. To characterise
this region of waveform space, we hold the electrode at $d=100\,\mu$m
(the cleanest distance from the strength--duration sweep) and grid four
parameters: anodic phase duration $T_p$, cathodic phase duration $T_n$,
anodic amplitude $A_p$ (with cathodic amplitude $A_n$ pinned by
$A_p T_p = A_n T_n$ to keep total injected charge zero), and pulse-train
frequency $f$. We record spike count, latency, and peak voltage per
configuration rather than searching for a threshold, so each grid point
is one forward simulation.

The grid has $T_p,\,T_n \in \{50,\,100,\,200,\,500\}\,\mu$s
(four values each), $A_p \in \{10,\,20,\,30,\,40,\,80\}\,\mu$A, and
pulse-train frequencies $f \in \{100,\,200,\,500,\,1000\}\,$Hz, for
$315$ valid charge-balanced forward simulations; the five invalid
combinations have $T_p + T_n$ equal to the $1000\,$Hz period. The full grid
runs in $16\,$min on the v5litepod-4 pod-slice via
\texttt{vmap}-over-configuration with hierarchical scan-checkpointing
(\cref{sec:res-throughput}). Spike crossings concentrate along the
asymmetric $T_p = 500\,\mu$s edge at low-to-mid amplitudes
($A_p \le 30\,\mu$A), and at $A_p = 80\,\mu$A for $T_p = 200\,\mu$s
(\cref{fig:biphasic-grid}); the symmetric diagonal $T_p = T_n$ is far
less effective at the same charge per pulse. The
implication for waveform design is direct: asymmetry is not a marginal
parameter --- on this cell, at this distance, swapping a symmetric
biphasic pulse for an asymmetric one with the same charge can be the
difference between sub-threshold and reliable spiking.

\begin{figure}[t]
  \centering
  \includegraphics[width=\linewidth]{figures/biphasic_grid.png}
  \caption{Charge-balanced biphasic grid on the BBP L2/3 Pyr cell at
  $d = 100\,\mu$m. Each panel is a $(T_p, T_n)$ heatmap of the maximum
  soma voltage over the four pulse-train frequencies, with one panel
  per anodic amplitude $A_p$. The dotted diagonal $T_p = T_n$ marks the
  symmetric-biphasic line; cells where $\mathit{vmax}$ crosses
  $0\,$mV correspond to a soma spike at some frequency in the sweep.
  White cells indicate configurations for which all simulated
  frequencies returned non-finite voltage traces and therefore no
  interpretable $\mathit{vmax}$; these occur mainly in the
  long-$T_p$ high-amplitude regime. Five period-filling configurations
  ($T_p = T_n = 500\,\mu$s at $1000\,$Hz) are absent from the
  forward-simulation set.}
  \label{fig:biphasic-grid}
\end{figure}

\subsection{Throughput scaling on TPU}\label{sec:res-throughput}

The point of the simulator extension is to support design-space
sweeps that were previously infeasible. We measure how it scales
across two regimes that probe different bottlenecks. A
\emph{cheap-cells} sweep isolates per-cell compile and kernel-dispatch
overhead at small problem size: \texttt{ncomp}$=2$ on the BBP Pyr
($\sim\!700$ compartments per cell after branchpoint expansion), $B$
up to $1024$ on the v5litepod-4 pod-slice. A \emph{full-fidelity}
sweep is the actual scientific workload: $50$-compartment branches on
the same Pyr cell ($\approx\!17\,500$ compartments per cell), $B$ up
to $4096$ on the same pod-slice. Each $B$ value is a single
\texttt{jax.jit(jax.vmap(\ldots))} dispatch; we report the cached
(JIT-warm) timing.

The cheap-cells curve scales linearly to $B=1024$ at $5.1\,$ms per cell
($\sim\!195$ cells/s); the full-fidelity curve scales linearly through
$B=1024$, peaks at $7.28$ cells per second at $B=2048$, and rolls over
at $B=4096$ as the four chips approach arithmetic-intensity saturation
(\cref{fig:throughput-scaling}). The full-fidelity regime depends on
the sparse-$G$ and \texttt{checkpoint\_lengths} choices in the Method
(\cref{sec:method}); without either, the $(T,\,B,\,N)$ voltage
history overflows the v5e's $16\,$GB chip memory at $B \ge 256$.

\begin{figure}[t]
  \centering
  \includegraphics[width=\linewidth]{figures/throughput_v5p4_scaling.png}
  \caption{Jaxley batch throughput on the v5litepod-4 pod-slice for two
  workload regimes: cheap cells ($700$ compartments per cell, blue) and
  full-fidelity BBP Pyr ($\sim\!17\,500$ compartments per cell, red).
  Throughput scales linearly through $B=1024$ in both regimes; the
  full-fidelity curve peaks at $7.28$ cells per second at $B=2048$ and
  saturates at $B=4096$. JIT-warm timings.}
  \label{fig:throughput-scaling}
\end{figure}

\paragraph{Comparison with NEURON.} The same $700$-compartment Pyr cell
on a $T_\text{sim} = 200\,$ms protocol runs at $0.286$ simulations per
second under serial single-core NEURON (the default BBP setup), against
$9.67$ simulations per second for Jaxley at $B = 1000$ on a single TPU
once the JIT cache is warm --- a $\sim\!34\times$ wall-clock speedup
(\cref{fig:throughput}).
NEURON has parallel-CPU variants such as
CoreNEURON~\cite{Kumbhar2019CoreNEURONA} and
experimental GPU support that we did not benchmark; the relevant
comparison for our purposes is between a CPU-serial reference and an
accelerator-batched formulation. AI accelerators are SIMT/SIMD
throughput devices, and a batched-\texttt{vmap} formulation is the
workload pattern they are designed for; a CPU-bound serial framework
cannot expose the same batch parallelism without a similar
reformulation. The practical consequence: the parameter sweeps in
\cref{sec:res-sd,sec:res-grid} finish in tens of minutes, putting
sweeps over the full waveform-and-geometry design space within reach
of a single TPU job.

\begin{figure}[t]
  \centering
  \includegraphics[width=\linewidth]{figures/bbp_throughput.png}
  \caption{Wall-clock comparison on the same $700$-compartment Pyr cell
  and $T_\text{sim} = 200\,$ms protocol: serial single-core NEURON
  (blue) versus Jaxley on a single TPU chip (red). At $B=1000$, NEURON
  takes $58.3\,$min (measured), Jaxley takes $1.7\,$min ($103.4\,$ms
  per simulation, JIT-warm), giving a $\sim\!34\times$ wall-clock
  speedup. The dashed blue line is the linear extrapolation of the
  per-NEURON-sim cost.}
  \label{fig:throughput}
\end{figure}

\subsection{Gradient-based design signal}\label{sec:res-grad}

For the differentiability we built into the pipeline to matter for
stimulation design, the gradients have to point in the right
direction --- and their direction has to mean something physical.
Establishing that gradients exist (\cref{sec:results-verify}) is a
prerequisite for, but not equivalent to, showing that they carry usable
design signal. The position gradient reported in
\cref{sec:results-verify} admits a direct physical reading. With
the electrode at $(0,\,100,\,0)\,\mu$m, the negative $y$-component of
$\partial v_\mathrm{soma,\,peak} / \partial \mathbf{x}_e$ means moving
the electrode farther from the cell along its dorsal axis reduces the
soma response, consistent with the $1/d$ decay in the field model; the
horizontal components measure the local field anisotropy induced by the
asymmetric placement of dendrites relative to the electrode. Free-form
optimisation of $w(t)$ against a spike-elicitation target under
charge-balance and amplitude constraints is a natural next step and is
left as future work (\cref{sec:limitations}). What this section
establishes is that the same pipeline that runs the empirical sweeps
of \cref{sec:res-sd,sec:res-grid} also returns design-relevant
gradients in a single backward pass --- placing direct waveform
optimisation, rather than gradient-free search, within reach of the
same compute budget.

\begin{figure}[t]
  \centering
  \includegraphics[width=0.78\linewidth]{figures/gradient_arrow.png}
  \caption{Position-gradient direction on the BBP L2/3 Pyr cell. Grey
  points are the $700$-compartment morphology projected onto the
  $y$-$z$ plane; the soma is marked with a black dot, the electrode
  ($\mathbf{x}_e = (0,\,100,\,0)\,\mu$m) with a black cross. The
  orange arrow shows the direction of
  $\partial v_\mathrm{soma,\,peak} / \partial \mathbf{x}_e$, i.e.\ the
  direction in which moving the electrode increases the soma peak.}
  \label{fig:gradient-arrow}
\end{figure}

% =============================================================
\section{Discussion}\label{sec:discussion}
% =============================================================

Genetic search over
waveforms~\cite{grill2015model,wongsarnpigoon2010energy} returns no
gradient direction for a candidate that is close to but short of an
activation target; surrogate networks trained on NEURON
outputs~\cite{hussain2024efficient} introduce an approximation error
that must be re-incurred whenever the underlying biophysical model
changes. Because the pipeline is JAX-traceable end-to-end, gradients
$\partial \ell / \partial w(t)$ of any differentiable scalar objective
$\ell$ with respect to the input waveform are returned in a single
backward pass and avoid both costs. As a
sharper receipt than the position gradient of \cref{sec:res-grad}, an
Adam loop on $\mathrm{MSE}(\hat{v}, v_\mathrm{target})$ on the full
BBP Pyr morphology recovers \texttt{gNaTs2T} from a $-15\,\%$
perturbation to within $+0.13\,\%$ in nine steps
(\texttt{fit\_demo}); the point of this fit is end-to-end gradient
flow on a realistic cell, not optimiser performance. A selectivity
Pareto sweep against a non-target cell, or a charge-minimising waveform
under a spike-elicitation constraint, follows the same template (a
differentiable spike-elicitation surrogate plus charge-balance and
amplitude constraints).

% =============================================================
\section{Limitations and future work}\label{sec:limitations}
% =============================================================

\paragraph{Limitations.}
\begin{itemize}
  \item \emph{Volume conductor.} The field model is a homogeneous,
  isotropic, infinite point-source monopole. Tissue anisotropy,
  finite electrode geometry (cuff, paddle, contact size), and the
  bone--CSF--grey-matter layering used in DBS planning are not
  represented.
  \item \emph{Quasi-static assumption.} Tissue capacitance is ignored;
  this is standard at neural-stimulation frequencies but breaks down at
  very short pulse edges.
  \item \emph{Float-32 caveat.} The production sweeps run in float-32
  after a $0.77\,\%$ threshold-difference check against float-64 on one
  HH cable configuration; the check should be repeated per cell type
  before reusing float-32 in a regime we have not validated.
  \item \emph{PV cell stability.} The full-morphology BBP L2/3 LBC
  basket cell is unstable in Jaxley's default integrator (NaN at
  $\sim\!30\,$ms even unstimulated), traced to BBP's
  $g_\mathrm{pas} = 1\!\times\!10^{-6}\,$S/cm$^2$ on PV apical and
  basal --- $30\times$ lower than the Pyr value and tolerated by NEURON
  but not by the Jaxley defaults. Single-compartment PV parity passes
  (\cref{sec:res-neuron}); all full-morphology runtime claims in this
  paper are therefore Pyr-only.
  \item \emph{Pyr Ih path-distance traversal.} The
  $-0.38\,$mV resting-potential offset against NEURON
  (\cref{sec:res-neuron}) is the residual of approximating BBP's
  path-distance apical $I_h$ gradient with Euclidean distance.
  Replacing the approximation with a parent-pointer breadth-first
  traversal would close the residual but we did not implement it.
  \item \emph{Solver floor.} \texttt{checkpoint\_lengths} bounds memory
  but not arithmetic. The backward-Euler scan dominates the per-step
  cost on a v5e chip and is what saturates the throughput curve at
  $B=2048$ in \cref{sec:res-throughput}.
\end{itemize}

\paragraph{Future work.}
Three lines of follow-up sit immediately on top of what is here. First,
restoring full PV cell coverage by stabilising the LBC integration ---
either with an explicit smaller-timestep regime, a leak-conductance
floor, or a stiff-aware solver --- and applying the BBP path-distance
$I_h$ gradient to PV apical and basal as in Pyr. Second, gradient-based
optimisation of the free-form waveform $w(t)$ against a
spike-elicitation surrogate, under charge-balance and amplitude
constraints, evaluated against the strength--duration baseline of
\cref{sec:res-sd}. Third, selectivity: two cell models with different
geometries placed in the same extracellular field, optimising $w(t)$
to drive a soma spike in the target cell while keeping the non-target
subthreshold. The Pareto frontier between activation and sparing is
exactly what the differentiable pipeline is built to trace; a
penalty-based formulation
$\mathcal{L}(w) = \lVert w \rVert_1 \Delta t
+ \lambda\,\tilde{F}(w;\theta_2)$ subject to
$\tilde{F}(w;\theta_1) \ge \tau$ keeps the problem smooth and lets a
$\lambda$-sweep trace the frontier directly.


% =============================================================
\section{Conclusion}\label{sec:conclusion}
% =============================================================

We have added extracellular stimulation to Jaxley by recovering the
discrete activating function from the simulator's own axial diffusion
operator and injecting it as an equivalent compartmental current through
the public stimulus API. The pipeline composes with \texttt{jit},
\texttt{vmap}, \texttt{grad}, and \texttt{jax.sharding} without a fork
of the simulator, agrees with Jaxley's internal voltage vector field to
$\sim\!10^{-13}$ relative error, and runs the BBP L2/3 pyramidal cell at
full $50$-compartment-per-branch fidelity on a four-chip v5e pod-slice.
What this enables, and what we leave to follow-up work, is gradient-based
search over the waveform itself --- the design problem that motivated
the extension.



\bibliographystyle{plain}
\bibliography{ref}

\appendix

% =============================================================
\section{Branchpoint elimination}\label{app:branchpoint-elimination}
% =============================================================

Jaxley's internal representation introduces branchpoint pseudo-nodes
at cable junctions to enforce Kirchhoff's current law. Partitioning
$G$ into real ($R$) and branchpoint ($B$) blocks gives
\[
  \begin{pmatrix} G_{RR} & G_{RB} \\ G_{BR} & G_{BB} \end{pmatrix}
  \begin{pmatrix} \mathbf{v}_R \\ \mathbf{v}_B \end{pmatrix}
  =
  \begin{pmatrix} [G\mathbf{v}]_R \\ [G\mathbf{v}]_B \end{pmatrix}.
\]
Branchpoint nodes carry no membrane, so $c_b A_b = 0$ at $b$ and the
cable equation at $b$ reduces to the algebraic constraint
$[G\mathbf{v}]_B = 0$, giving $\mathbf{v}_B = -G_{BB}^{-1}G_{BR}\mathbf{v}_R$
and the Schur complement
\[
  \tilde{G} \;=\; G_{RR} - G_{RB}\,G_{BB}^{-1}\,G_{BR}.
\]
The same operator applies to the extracellular forcing:
$\tilde{G}\boldsymbol{\phi}_R$ gives the activating function at the
real compartments, where $\boldsymbol{\phi}_R$ is the field sampled at
real compartment centres only. The reduced operator acts only on real
compartments, so no $\boldsymbol{\phi}_B$ value enters the forcing. Our implementation delegates this assembly to a
private Jaxley helper rather than reconstructing it from scratch, so
the operator stays in sync with upstream solver changes.


% =============================================================
\section{Equivalent-current encoding}\label{app:equivalent-current}
% =============================================================

The activating function $G\boldsymbol{\phi}$ has units mV/ms (a
voltage rate), but Jaxley's stimulus API accepts current in nA.
Internally Jaxley converts
$i_{\text{ext}}\,[\mu\text{A/cm}^2] = i_{\text{nA}}\,[\text{nA}] /
A\,[\mu\text{m}^2]\cdot 10^5$ and adds $i_{\text{ext}}/c$ to $dv/dt$.
To inject $[G\boldsymbol{\phi}]_j$ on the other side of this
conversion, we encode
\[
  i_{\text{ecs},j}\,[\text{nA}] \;=\;
    c_j \cdot [G\boldsymbol{\phi}(t)]_j \cdot \frac{A_j}{10^5}.
\]
The $c_j$ and $A_j$ factors cancel in the round trip; they are an
encoding artifact of using the public API rather than modifying the
solver. In matrix form,
\[
  \mathbf{I}_{\text{ecs}} \;=\;
    \left(\mathbf{c} \odot \frac{\mathbf{A}}{10^5}\right)
    \odot (G\boldsymbol{\Phi})
    \;\in\; \mathbb{R}^{N\times T}.
\]


% =============================================================
\section{BBP cell-model implementation notes}\label{app:bbp-implementation}
% =============================================================

Two implementation issues affect reproducibility of the BBP cell models
in Jaxley but are unrelated to the paper's results.

\paragraph{SWC soma conversion.} The BBP source morphologies represent
the soma as a two-dimensional contour, which the SWC format cannot
encode. The Neurolucida-to-SWC conversion tool emits soma points with
$r = 0$, which Jaxley turns into degenerate zero-area compartments
that produce NaN voltages at the first integration step. We patch the
converted SWC files at load time so the soma points carry a
non-degenerate radius.

\paragraph{Calcium-channel insertion order.} Calcium signalling is
sensitive to channel insertion order. The calcium Nernst reversal
must be established before any calcium-carrying channel, which in
turn must precede the calcium pump and the calcium-activated
potassium channel.
\texttt{cell\_factory.\_insert\_calcium\_chain} enforces
\texttt{CaNernstReversal} $\to$ \texttt{CaHVA/CaLVA} $\to$
\texttt{CaPump} $\to$ \texttt{SKE2}; out-of-order insertion produces
uninitialised state references at integration time.


% =============================================================
\section{Apical \texorpdfstring{$I_h$}{Ih} distance-dependent
gradient}\label{app:ih-gradient}
% =============================================================

The \texttt{cADpyr229} apical $I_h$ conductance follows BBP's
distance-dependent gradient
\[
  \bar{g}_{I_h}(d)
    = \bigl(-0.8696 + 2.087\,e^{0.0031\,d}\bigr)
      \cdot 8\!\times\!10^{-5}\;\mathrm{S/cm}^2,
\]
where $d$ is path distance from the soma centre in $\mu$m. We
approximate $d$ with the Euclidean distance to the soma centre, which
is close for roughly-vertical apical trees. Before the fix the
resting-potential offset against NEURON was $-0.62\,$mV; the
Euclidean approximation closes it to $-0.38\,$mV, and a path-distance
traversal would close it further (\cref{sec:limitations}).


% =============================================================
\section{Reproducibility}\label{app:reproducibility}
% =============================================================

The build pipeline binds the entire compute environment, not just the
Python packages. A \texttt{requirements.txt} inside a Dockerfile pins
Python packages and a base OS image but leaves the C and
C\nolinkurl{++} compilers, the CUDA toolchain, the Python interpreter
build, and the cloud topology unspecified --- each of which can drift
between machines and change numerical results. The stack we use
closes those gaps.

\paragraph{Python dependencies (\texttt{uv}).} A Rust-implemented
package manager that resolves direct and transitive Python
dependencies and records every resulting wheel and its hash in a
lockfile (\texttt{uv.lock}). Two machines sharing the lockfile
install the same versions of every Python package, regardless of
install order.

\paragraph{System environment (Nix).} A Nix flake (\texttt{flake.nix})
specifies the Python interpreter (3.11), CUDA runtime, JAX, the C
compiler that built them, and command-line tooling. Store paths are
determined by the hash of the full transitive build graph (pinned via
\texttt{flake.lock}), so the same flake on the same target system
produces the same dependency closure.

\paragraph{Cloud infrastructure (OpenTofu).} The TPU pod-slice, Cloud
SQL Postgres database, and GCS artifact bucket are declared as
OpenTofu configuration in \texttt{infra/tofu/}; \texttt{tofu apply}
brings the topology up, \texttt{tofu destroy} tears it down.

\paragraph{Reproducing the figures.} A reader entering the Nix shell
gets a deterministic Python and CUDA environment; the project
\texttt{README} carries the exact commands. Reproducing the figures
from pinned data does not require any TPU-scale workload.

\paragraph{Artifact set.} The curated data backing every figure ---
parity \texttt{.npz} files, the strength--duration and biphasic-grid
Zarr stores, the throughput sweeps, and the gradient receipt ---
totals a few megabytes and is published as a separate artifact
distribution rather than committed to the source repository, available
at \url{https://github.com/Marco-Christiani/jaxley-extracellular/releases/tag/paper-v0.1}.


\end{document}

% so a single source can produce both targets.
\providecommand{\iflong}{}
\ifx\iflong\empty
  \newif\iflong
  \longtrue
\fi

% Stand-in markers. Anything tagged \todo{...} is a placeholder for content
% that has been sketched but not written, or where the author still needs to
% confirm a number, citation, or figure. Render in the margin so the gaps are
% visible at a glance during review.
\newcommand{\todo}[1]{\textcolor{red}{\textbf{[TODO: #1]}}}
\newcommand{\sketch}[1]{\textcolor{gray}{\textit{(sketch: #1)}}}
\newcommand{\fig}[1]{\textcolor{blue}{\textbf{[FIG: #1]}}}

\title{Differentiable Extracellular Stimulation
of Compartmental Neuron Models}
\author{%
  Marco Christiani \\
  Carnegie Mellon University \\
  \texttt{marcochr@andrew.cmu.edu}
}

\graphicspath{{./}{figures/}}

\begin{document}
\maketitle

\begin{abstract}
We extend Jaxley, a JAX-native compartmental neuron simulator, with
extracellular stimulation. The extension constructs the discrete activating
function from Jaxley's pre-normalized axial diffusion operator and injects it
as an equivalent compartmental current through the public stimulus API; no
fork or solver patch is required. Every step is JAX-traceable, so
\texttt{jit}, \texttt{vmap}, \texttt{grad}, and \texttt{jax.sharding} compose
with the pipeline without modification. We verify the operator against
Jaxley's internal voltage vector field to $\sim\!10^{-13}$ relative error in
float64 and recover $O(\Delta x^2)$ convergence to the analytical activating
function on a uniform straight cable. We then exercise the pipeline on
Hodgkin--Huxley cables and on the full-morphology Blue Brain Project
L2/3 pyramidal cell, with the parvalbumin basket cell included only at
the single-compartment channel-kinetics level since its full
morphology does not run stably in our setup. We validate against
NEURON in a three-step cross-check that reaches a median per-site
voltage RMSE of $0.446\,$mV on the full-morphology Pyr cell under
extracellular stimulation, sweep stimulation parameters across pulse
width, polarity, distance, and frequency on a four-chip TPU pod-slice,
and report a $\sim\!34\times$ wall-clock speedup over serial
single-core NEURON on the same Pyr cell. Gradients
of voltage observables flow through the full simulator with respect to
electrode position, waveform parameters, and tissue conductivity,
placing direct waveform optimisation within reach of the same compute
budget.
\end{abstract}


% =============================================================
\section{Introduction}\label{sec:intro}
% =============================================================

Clinical neurostimulation devices, including deep brain stimulators and
spinal cord stimulators, deliver pulsed extracellular currents through
electrodes placed near the target tissue. Therapeutic outcome depends on a
small set of waveform parameters --- pulse width, shape, amplitude,
frequency, polarity --- whose effect on neural activation is nonlinear and
sensitive to local geometry. In present clinical practice these parameters
are tuned by hand, guided by symptom observation across patient visits.
Computational design has been explored as an
alternative, but the simulators in widespread use are slow when the goal is
to sweep the design space.

The standard pipeline pairs a volume-conductor field model with the NEURON
simulator~\cite{Hines2001NeuronAT} to integrate the resulting
transmembrane response. The default BBP NEURON workflow runs serially
on a single CPU core, so existing studies that search for energy-efficient or selective
waveforms have either used genetic
algorithms~\cite{grill2015model,wongsarnpigoon2010energy} or trained
surrogate neural networks on NEURON
outputs~\cite{melissano2024efficient}. Genetic search returns no gradient
and requires many forward simulations per design. Surrogates are fast at
inference but introduce approximation error and have to be retrained
whenever the underlying biophysical model changes.

Jaxley~\cite{deistler2025jaxley} is a recently published JAX-native
compartmental simulator. It supports automatic differentiation, JIT
compilation, \texttt{vmap}, and the modern \texttt{jax.sharding} API for
multi-device execution. It does not currently expose extracellular
stimulation. We add it without forking the simulator: the extension
constructs the discrete activating function from Jaxley's own axial
diffusion operator and injects the result as a compartmental current
through the public stimulus API.

\paragraph{Contributions.}
\begin{enumerate}
  \item A differentiable extracellular-stimulation pipeline through
  Jaxley's public stimulus API, with no fork of the upstream simulator.
  \item A derivation of the discrete activating function
  $G\boldsymbol{\phi}$ from Jaxley's pre-normalised voltage diffusion
  operator $G$.
  \item Numerical verification: agreement with Jaxley's internal voltage
  vector field to $\sim\!10^{-13}$ relative error in float64,
  $O(\Delta x^2)$ convergence to the analytical activating function, and
  end-to-end gradient flow through \texttt{jx.integrate}.
  \item Three-step cross-validation against NEURON --- single-compartment
  kinetics, full-morphology Pyr intracellular, and full-morphology Pyr
  under extracellular stimulation --- with characterised residuals at
  each step.
  \item TPU-scale empirical results on the BBP L2/3 pyramidal cell:
  a strength--duration sweep across distance and polarity, a
  charge-balanced biphasic grid of $315$ valid configurations, and a
  throughput--scaling curve at $\sim\!34\times$ wall-clock speedup over
  serial NEURON at $B = 1000$.
\end{enumerate}


% =============================================================
\section{Background and related work}\label{sec:background}
% =============================================================

\subsection{Extracellular potential from point-source electrodes}

At neural stimulation frequencies, displacement currents in tissue are
negligible compared with conduction currents, and Maxwell's equations
reduce to a quasi-static formulation in which the extracellular potential
satisfies Laplace's equation away from sources. Modelling each electrode
contact as a point source in an infinite, homogeneous, isotropic
conductor of conductivity $\sigma$ gives the closed-form monopole
$\phi(\mathbf{x}) = I / (4\pi\sigma\,d)$, with $d$ the source-to-field
distance~\cite{10.1093/acprof:oso/9780195058239.001.0001,10.1093/acprof:oso/9780195050387.001.0001}.
Multiple electrodes follow by superposition. We use this model throughout; the homogeneity and isotropy
assumptions are limitations we return to in \cref{sec:limitations}.

\subsection{The activating function}

For a uniform straight cable in an extracellular field $\phi(x,t)$,
Rattay~\cite{rattay1986analysis} showed that the driving term for
transmembrane voltage is proportional to
$\partial^2 \phi / \partial x^2$ along the cable; McNeal's earlier
compartmental treatment of myelinated axons under monopolar
stimulation~\cite{Mcneal1976AnalysisOA} contains the same idea in
discrete form. On an arbitrary compartmental tree the discrete
analogue is $[G\boldsymbol{\phi}]_j$, where $G$ is the axial diffusion
operator that the cable solver already maintains; the continuous second
derivative is the special case of a uniform unbranched cable. We exploit
this in \cref{sec:method}: rather than recomputing
$\partial^2 \phi / \partial x^2$ along each branch, we apply Jaxley's
$G$ directly to the sampled extracellular potential.

\subsection{Differentiable neuron simulators}

Jaxley~\cite{deistler2025jaxley} is the only published JAX-native
compartmental simulator we are aware of that supports HH-style channel
kinetics, branched morphologies, and reverse-mode autodiff through the
time integrator. Diffrax~\cite{DBLP:journals/corr/abs-2202-02435} provides
differentiable ODE solvers in JAX, but it does not model compartmental
neurons. Brian2~\cite{Stimberg2019} is
Python-native and considerably faster than NEURON, but it is not built
for automatic differentiation. We reuse Jaxley's solver through its public
stimulus API rather than forking the simulator: the extracellular extension
is a JAX module that composes with the rest of the framework without
modifying its internals.


% =============================================================
\section{Method}\label{sec:method}
% =============================================================

We adopt the notation of \cref{tab:notation}. The pipeline is:
\begingroup
\setlength{\abovedisplayskip}{4pt}
\setlength{\belowdisplayskip}{4pt}
\begin{equation}
  I_e(t) \xrightarrow{\text{field}} \boldsymbol{\Phi}
  \xrightarrow{G\,\cdot} \mathbf{f}
  \xrightarrow{\text{encode}} \mathbf{I}_{\text{ecs}}
  \xrightarrow{\text{integrate}} \mathbf{v}(t).
  \label{eq:pipeline}
\end{equation}
\endgroup

\begin{table}[t]
\centering
\small
\caption{Notation. Lengths in $\mu$m, currents in $\mu$A or nA, voltages
in mV, time in ms, conductivity in S/m.}
\label{tab:notation}
\begin{tabular}{@{}lll@{}}
\toprule
Symbol & Meaning & Units \\ \midrule
$N$ & Number of compartments & --- \\
$T$ & Number of timesteps & --- \\
$\mathbf{x}_j,\,\mathbf{x}_e$ & Compartment center, electrode position & $\mu$m \\
$d_{je}$ & $\lVert \mathbf{x}_j - \mathbf{x}_e \rVert$ & $\mu$m \\
$I_e(t)$ & Electrode current (cathodic negative) & $\mu$A \\
$\sigma$ & Extracellular conductivity & S/m \\
$\phi_j(t),\,\boldsymbol{\Phi}$ & Extracellular potential, $\boldsymbol{\Phi}_{jt}=\phi_j(t)$ & mV \\
$v_j(t),\,\mathbf{v}$ & Transmembrane voltage & mV \\
$G$ & Axial voltage diffusion operator & $1/$ms \\
$c_j,\,A_j$ & Specific capacitance, surface area & $\mu$F/cm$^2$,\,$\mu$m$^2$ \\
$f_j(t) = [G\boldsymbol{\phi}(t)]_j$ & Activating function at compartment $j$ & mV/ms \\
\bottomrule
\end{tabular}
\end{table}

\paragraph{Extracellular potential from point sources.}
For an electrode $e$ delivering $I_e(t)$ in a homogeneous, isotropic,
infinite volume conductor of conductivity $\sigma$, superposition of
monopoles gives
\[
  \phi_j(t) \;=\; \sum_e \frac{I_e(t)}{4\pi\sigma\,d_{je}}.
\]
Defining the transfer factor $h_{je} = 10^3 / (4\pi\sigma\,d_{je})$
(in $\mathrm{mV}/\mu\mathrm{A}$, with the $10^3$ absorbing the unit
conversion from $\mu$A and $\mu$m to mV) factorises space and time, so
that for a single electrode
$\boldsymbol{\Phi} = \mathbf{h}\,\mathbf{I}^\top \in \mathbb{R}^{N\times T}$.

\paragraph{The discrete activating function.}
Jaxley solves the spatially discretised cable equation
\[
  c_j A_j \frac{dv_j}{dt}
  = \sum_{k\sim j} g_{jk}(v_k - v_j) - A_j I_{\text{ion}}(v_j,\mathbf{s}_j) + i_{\text{ext},j},
\]
where $g_{jk}$ is the raw axial conductance (units S) between
compartments $j$ and $k$, $c_j$ is specific membrane capacitance, and
$A_j$ is compartment surface area. Internally, Jaxley folds $c_j A_j$
and unit conversions into a single voltage diffusion operator $G$ with
entries $\tilde g_{jk} = g_{jk} / (c_j A_j)$ (in $1/$ms; note the
asymmetric row-$j$ normalisation, which is why $G$ is symmetric only
when $c_j A_j$ is uniform across compartments), so that
$\dot{\mathbf{v}} = G\mathbf{v} + (\text{membrane terms})$. The sign
pattern of $G$ is the weighted graph Laplacian:
$G_{jk} = +\tilde g_{jk}$ for adjacent $k\neq j$,
$G_{jj} = -\sum_{k\neq j} \tilde g_{jk}$, and zero otherwise. Rows sum
to zero (current conservation).

An extracellular field $\phi_j(t)$ shifts the axial driving force from
$v_k - v_j$ to $(v_k+\phi_k)-(v_j+\phi_j)$; by linearity of $G$, this
is equivalent to adding $G\boldsymbol{\phi}$ as a forcing term, the
discrete analogue of Rattay's continuous activating function
$\rho_a^{-1}\partial^2\phi/\partial x^2$~\cite{rattay1986analysis}.
On a uniform straight cable with constant radius and compartment
length $\Delta x$, $[G\boldsymbol{\phi}]_j$ at interior compartments
reduces to a constant multiple of the standard
$[1,-2,1]/\Delta x^2$ stencil for $\partial^2\phi/\partial x^2$, with
$O(\Delta x^2)$ truncation error (verified in
\cref{sec:results-verify}).
On branched morphologies, we delegate the elimination of the
branchpoint pseudo-nodes that Jaxley introduces at cable junctions to
its own solver routine, keeping the operator in sync with upstream
solver changes (\cref{app:branchpoint-elimination}).

\paragraph{Equivalent injected current.}
The activating function $G\boldsymbol{\phi}$ has units mV/ms (a voltage
rate), but Jaxley's stimulus API accepts current in nA. We encode the
forcing as a per-compartment current $\mathbf{I}_{\text{ecs}}$ chosen
so that, after Jaxley's internal current-to-voltage-rate conversion,
the integrator sees exactly $G\boldsymbol{\phi}$. The unit-conversion
algebra is given in \cref{app:equivalent-current}; the membrane
capacitance and area factors that appear in the encoding cancel
exactly in the round trip and are an artifact of working through the
public API rather than modifying the solver.

\paragraph{Implementation.}
The pipeline of \cref{eq:pipeline} is realised as four small JAX-native
modules wired together by an experiment object.
\texttt{field.point\_source\_potential} computes $\boldsymbol{\Phi}$ from
electrode positions, currents, and $\sigma$, with a $1\,\mu$m floor on
$d_{je}$ to regularise the $1/d$ singularity when an electrode lands inside
a compartment.
\texttt{discretization.build\_voltage\_operator\_G} returns $G$ by
delegating to Jaxley's \texttt{\_compute\_transition\_matrix} (the same
routine the simulator uses for its backward-Euler step) and stripping
branchpoint pseudo-nodes.
\texttt{equivalent\_current.phi\_e\_to\_ecs\_nA} applies the encoding
above.
An \texttt{ECSExperiment} class precomputes $G$, $\mathbf{c}$, $\mathbf{A}$,
and the compartment centres in its constructor, and exposes
\texttt{simulate\_waveform} for a single integration and
\texttt{simulate\_and\_extract} for one integration plus response-feature
extraction.

All operations are JAX-traceable: \texttt{jit}, \texttt{vmap},
\texttt{grad}, and \texttt{jax.sharding.NamedSharding} compose with
the pipeline without modification, so a one-dimensional device mesh is
enough to data-parallelise the batch axis across a TPU pod-slice. Two
implementation choices matter at scale and unlock the full-fidelity
throughput regime of \cref{sec:res-throughput}: $G$ is stored as a
sparse \texttt{BCOO} matrix to avoid the $(B,\,N,\,N)$ XLA broadcast
that dense $G$ triggers under \texttt{vmap}, and
\texttt{simulate\_waveform} forwards a \texttt{checkpoint\_lengths}
factorisation of $T$ into \texttt{jx.integrate} so the backward-Euler
scan rematerialises in $\sim\!\sqrt{T}$ chunks rather than holding the
full $(T,\,B,\,N)$ voltage history in HBM.


% =============================================================
\section{Cell models}\label{sec:cells}
% =============================================================

\subsection{Hodgkin--Huxley straight cable}\label{sec:cells-hh}

The verification studies of \cref{sec:results-verify} and the pilot
strength--duration sweep of \cref{sec:res-sd} use a uniform
Hodgkin--Huxley cable. Defaults reproduce the configuration set by
\texttt{experiment.make\_hh\_cable\_experiment}: $50$ compartments
along a $1.25\,$mm cable, fibre radius $10\,\mu$m, axial resistivity
$\rho_a = 100\,\Omega\cdot$cm, with the standard Jaxley HH channel set on
every compartment. A single electrode sits $50\,\mu$m above the first
compartment in a medium of conductivity $\sigma = 0.3\,$S/m, the typical
value for cortical grey matter.

\subsection{BBP L2/3 cells (cADpyr229, cNAC187)}\label{sec:cells-bbp}

For morphologically detailed simulations we adapt two L2/3 cell models
from the Blue Brain Project~\cite{markram2015reconstruction}: the
regular-spiking pyramidal cell
\texttt{L23\_PC\_cADpyr229} and the parvalbumin-expressing basket cell
\texttt{L23\_LBC\_cNAC187}. Channel kinetics are reused from
\texttt{jaxley\_mech.channels.l5pc}, which provides Jaxley-compatible
implementations of the BBP L5 channel set; conductances and passive
parameters are read from the BBP biophysics specification and applied per
section group (soma, axon, apical, basal).

Two implementation details affect reproducibility but do not bear on
the paper's results: the SWC conversion of the BBP soma contour, and
the calcium-channel insertion order in Jaxley. Both are documented in
\cref{app:bbp-implementation}.

A third detail is needed for NEURON parity. The \texttt{cADpyr229}
apical $I_h$ conductance follows a distance-dependent gradient
specified in BBP's biophysics file with a path-distance argument; we
apply it under a Euclidean-distance approximation to the soma centre,
which closes a $0.6\,$mV resting-potential offset that a uniform
somatic value otherwise leaves (closed-form gradient and
implementation in \cref{app:ih-gradient}). Long SWC branches are
subdivided at load time (\texttt{max\_branch\_len}$\,=\,100\,\mu$m) so
the compartment count tracks what NEURON's \texttt{geom\_nseg} routine
produces for the same morphology; we found this necessary for
spike-timing parity (\cref{sec:discussion}).


% =============================================================
\section{Numerical verification}\label{sec:results-verify}
% =============================================================

Three numerical checks anchor the rest of the paper: that $G$ matches what
Jaxley's solver computes internally, that $G\boldsymbol{\phi}$ recovers the
continuous activating function as the discretisation refines, and that
gradients propagate through the full simulator. All three are exercised in
the test suite (\texttt{test\_ecs\_operator\_consistency} and
\texttt{test\_ecs\_grad}).

\paragraph{Operator round-trip.}
We compare $G\mathbf{v}$ against Jaxley's internal axial term on a
50-compartment cable in float64, isolating the axial term by zeroing
the membrane contributions (\texttt{voltage\_terms}\,$=\,0$ and
\texttt{constant\_terms}\,$=\,0$). The two outputs agree to
$\sim\!10^{-13}$ per-compartment relative error
(\cref{fig:verification}, left); once $G\mathbf{v}$ matches Jaxley's
own axial term, $G\boldsymbol{\phi}$ is the correct extracellular
forcing on the same tree.

\paragraph{$O(\Delta x^2)$ convergence.}
For a uniform straight cable with a point-source electrode the activating
function has the closed form $\rho_a^{-1}\,\partial^2\phi/\partial x^2$.
We compare $G\boldsymbol{\phi}$ against this analytical reference on the
central $10\,\%$ of the cable (boundary effects dominate near the ends)
across compartment counts $\{50,\,100,\,200,\,400\}$. The relative error
decreases by a factor of $\sim 4$ per doubling
($3.95\times,\,3.99\times,\,4.00\times$), consistent with the
$O(\Delta x^2)$ rate of the underlying $[1,-2,1]/\Delta x^2$ stencil
(\cref{fig:verification}, right).

\begin{figure}[t]
  \centering
  \includegraphics[width=\linewidth]{figures/verification.png}
  \caption{Numerical verification of the discrete activating function.
  Left: histogram of per-compartment relative error of $G\mathbf{v}$ vs
  Jaxley's internal \texttt{\_voltage\_vectorfield} on a $50$-compartment
  HH cable in float64; the median is $\sim\!3\!\times\!10^{-16}$ and no
  compartment exceeds $\sim\!10^{-13}$. Right: relative error of
  $G\boldsymbol{\phi}$ versus the closed-form
  $\rho_a^{-1}\partial^2\phi/\partial x^2$ on the central $10\,\%$ of the
  cable, plotted against $\Delta x$ on log--log axes; the dashed line is
  a reference slope-2, and the box reports the per-doubling error
  ratios.}
  \label{fig:verification}
\end{figure}

\paragraph{Gradient flow.}
The pipeline is JAX-traceable end to end, so \texttt{jax.grad} of any
voltage observable propagates through the field model, the operator,
the encoding, and Jaxley's integrator. Unit tests confirm finite,
non-zero gradients with respect to electrode position, waveform, and
tissue conductivity. At full BBP scale, the position gradient through
a $700$-compartment Pyr cell under a $-50\,\mu$A, $1\,$ms cathodic
pulse from electrode $(0, 100, 0)\,\mu$m is
\[
  \frac{\partial v_{\mathrm{soma,\,peak}}}{\partial \mathbf{x}_e}
  = (\,0.330,\,-0.371,\,-0.243\,)\;\mathrm{mV}/\mu\mathrm{m},
\]
returned in $37.6\,$s on a v5e single chip. The full receipt is part
of the artifact set described in the reproducibility appendix.


% =============================================================
\section{Empirical results}\label{sec:results-empirical}
% =============================================================

This section answers four questions that a stimulation designer or
methodologist would put to the simulator. Does it agree with NEURON on
a realistic morphology before we trust any extracellular result? At a
given electrode distance and pulse width, what amplitude is needed to
evoke a spike? Where in the asymmetric biphasic waveform space does
the cell respond, given the same charge per pulse? And does the
infrastructure scale to design-space sweeps that NEURON cannot run? We
close with a gradient receipt that motivates the next step --- direct
optimisation of the waveform itself. All runs reported below are part
of the curated artifact set described in the reproducibility appendix;
wall-clock figures are on a single v5e chip unless stated otherwise.

\subsection{Validation against NEURON}\label{sec:res-neuron}

Before any stimulation result from this pipeline is trustworthy, the
simulator has to agree with NEURON on the same cell. We validate in
three steps, each adding one new mechanism on top of the previous so
that residuals can be attributed.

\paragraph{Setup and confounders.}
The NEURON reference uses the same BBP cADpyr229 morphology that
Jaxley loads, with three confounders equalised before measurement.
(i) Axon myelination is enabled by default in the BBP NEURON template;
we disable it for the comparison since our Jaxley implementation does
not model myelin. (ii) The \texttt{jaxley\_mech.channels.l5pc} ports
apply a fixed $Q_{10}$ temperature correction; we run NEURON at
$34\,^\circ$C to match. (iii) NEURON's \texttt{geom\_nseg} produces $1131$
segments on this morphology, while Jaxley with \texttt{ncomp}$\,=\,2$
and \texttt{max\_branch\_len}$\,=\,100$ produces $700$ compartments;
we report the mismatch rather than forcing them to match.

\paragraph{Step 1: channel kinetics.} We reduce the morphology to a
$20\,\mu$m$\times 20\,\mu$m cylinder (\texttt{parity\_single\_soma}),
so that the comparison exercises only channel kinetics and ion
reversals. With the full BBP Pyr and PV channel sets installed, spike
counts agree to $\pm 1$ over $\sim\!50$ spikes per cell (Pyr: $18$
NEURON / $16$ Jaxley; PV: $50$/$51$), and AP peak heights agree to
$\pm 0.04\,$mV across both cells. The match is consistent with a
source-level diff of the relevant \texttt{jaxley\_mech} channels
against the BBP \texttt{.mod} files.

\paragraph{Step 2: full-morphology Pyr, intracellular.} We instantiate
the full cADpyr229 morphology in both simulators and inject a
$0.5\,$nA somatic step current with extracellular coupling held out
(\texttt{parity\_bbp\_pyr}). Both simulators emit $3$ spikes
(\cref{fig:parity-intracellular}); first-spike latency differs by
$-50\,\mu$s ($2$ integration steps at $\Delta t = 0.025\,$ms), and AP
peak by $+0.49\,$mV. Subthreshold
waveform RMSE is $0.275\,$mV with Pearson $r = 0.9995$;
suprathreshold full-trace RMSE is $6.032\,$mV with $r = 0.9253$, where
the latter is dominated by spike-edge timing offsets. The
$-0.38\,$mV resting-potential residual is explained by our
Euclidean-distance approximation to BBP's path-distance apical $I_h$
gradient (\cref{sec:cells-bbp}); a path-distance traversal would close
it further.

\paragraph{Step 3: full-morphology Pyr, extracellular.} The headline
claim of this paper is that ECS works on a realistic branched cell.
We test it by comparing voltage traces with NEURON when both
simulators see the same point-source field on the same Pyr
morphology. After aligning Jaxley's compartment centres and NEURON's
segment readouts to the same physical sites
(\texttt{parity\_ecs.bbp\_pyr}), per-site voltage RMSE is
$0.387,\,0.553,\,0.505,\,0.322\,$mV at soma, apical, basal and axon
respectively (median $0.446\,$mV); the site-local $\phi$ matches
between the two simulators (\cref{fig:parity-ecs}). The ECS residuals should be read against
the Step-2 baseline, since the underlying solver and discretisation
differences carry over. Combined with the analytical
$O(\Delta x^2)$ convergence of \cref{sec:results-verify} on a uniform
straight cable, this anchors the activating-function pipeline at both
geometric extremes: closed-form correctness on a straight cable and
external parity on a branched real cell. The practical consequence is
that ECS-on-Pyr in this pipeline reproduces NEURON behaviour to
sub-millivolt residuals on a realistic morphology --- enough fidelity
to support design-space sweeps that NEURON cannot run in the same time
budget (\cref{sec:res-throughput}).

\begin{figure}[t]
  \centering
  \includegraphics[width=\linewidth]{figures/bbp_intracellular_parity.png}
  \caption{Step 2 (full-morphology Pyr, intracellular) cross-check.
  NEURON (black) and Jaxley (dashed) soma voltage traces under $0.1\,$nA
  subthreshold and $0.5\,$nA suprathreshold step currents on the same
  cADpyr229 morphology. Subthreshold RMSE $0.275\,$mV
  ($r = 0.9995$); suprathreshold RMSE $6.032\,$mV
  ($r = 0.9253$), dominated by spike-edge timing offsets
  attributable to the $700$-vs-$1131$ compartmentalisation difference.}
  \label{fig:parity-intracellular}
\end{figure}

\begin{figure}[t]
  \centering
  \includegraphics[width=\linewidth]{figures/bbp_ecs_parity.png}
  \caption{Step 3 (full-morphology Pyr, extracellular) cross-check.
  NEURON (black) and Jaxley (dashed) voltage traces at the soma and
  apical readout sites after aligning the two simulators' compartment
  centres to within $0\,\mu$m on the same Pyr morphology, plus per-site
  RMSE bars and the site-local $\phi$ that drives the ECS coupling.
  Median per-site RMSE $0.446\,$mV; the residuals carry over from the
  Step-2 baseline rather than reflecting a new ECS-specific
  discrepancy.}
  \label{fig:parity-ecs}
\end{figure}

\subsection{Strength--duration on the HH cable}\label{sec:res-sd}

At a given electrode distance and pulse width, what amplitude does it
take to evoke a spike? The strength--duration curve is the standard
first answer to that question, and underlies almost every clinical and
engineering decision downstream of it~\cite{4122088}. We define the threshold $A^*$ for a fixed
configuration $(d, \text{polarity}, t_p)$ as the smallest electrode
amplitude that elicits a single soma spike in the simulation:
\[
  A^*(d,\,\text{polarity},\,t_p)
  \;=\; \inf\bigl\{A \ge 0\,:\,F(w_{A,t_p};\,d) \text{ spikes}\bigr\}.
\]
Each $A^*$ is found by a one-dimensional binary search in $A$, with the
spike-detection criterion threshold-crossing of the soma voltage (see
\texttt{response.detect\_spike}). The full sweep is one binary search per
$(d,\,\text{polarity},\,t_p)$ triple --- four electrode distances
$\{25,\,50,\,100,\,200\}\,\mu$m, three waveform polarities (monophasic
cathodic, monophasic anodic, charge-balanced biphasic), and nine pulse
widths from $0.05$ to $2\,$ms --- for $108$ thresholds in total.

The search uses a uniform $[0,\,5000]\,\mu$A bracket with $14$
bisection iterations, giving a resolution of $\sim\!0.3\,\mu$A; spike
detection stays stable across this range for the HH cable. The whole
sweep runs in float-32. Cathodic and biphasic thresholds are well-resolved at
every $(d,\,t_p)$ pair and follow the textbook strength--duration
shape, declining toward a rheobase as $t_p$ grows
(\cref{fig:strength-duration}). Monophasic anodic elicits spikes only
at short pulse widths and short distances; it saturates at the bracket
ceiling for longer pulses. This polarity asymmetry is the discrete
analogue of the activating-function sign argument: cathodic stimulation
sets up a depolarising forcing in the central segments where the cable
fires first, whereas anodic stimulation depolarises the cable ends,
which only triggers under brief, geometrically-favourable
conditions~\cite{rattay1986analysis,Mcneal1976AnalysisOA}. For a
designer choosing parameters, the takeaway is that
ms-scale cathodic pulses minimise charge per spike at every distance
in the sweep, charge-balanced biphasic costs roughly two-fold more,
and pure anodic stimulation is only viable in the short-pulse corner.

\begin{figure}[t]
  \centering
  \includegraphics[width=\linewidth]{figures/strength_duration.png}
  \caption{Strength--duration curves on the $50$-compartment HH cable.
  One panel per electrode distance; one curve per waveform polarity.
  Threshold amplitude $|A^*|$ falls toward a rheobase as pulse width
  grows, and grows with electrode distance as the $1/d$ field decay
  implies. The dotted line is the binary-search ceiling ($5000\,\mu$A);
  pulse-width points missing from a curve indicate no spike below that
  ceiling, which characterises the anodic regime at long widths.}
  \label{fig:strength-duration}
\end{figure}

\subsection{Charge-balanced biphasic grid}\label{sec:res-grid}

Asymmetric biphasic pulses are the clinical lever for trading
activation selectivity against charge efficiency: at the same total
delivered charge per pulse, where on the $(T_p, T_n)$ plane does the
cell actually respond? A symmetric charge-balanced biphasic pulse has
two equal-and-opposite phases of equal duration; a clinical waveform
almost always has unequal phase durations (longer-low cathodic +
short-high anodic, or vice versa)~\cite{wongsarnpigoon2010energy}. To characterise
this region of waveform space, we hold the electrode at $d=100\,\mu$m
(the cleanest distance from the strength--duration sweep) and grid four
parameters: anodic phase duration $T_p$, cathodic phase duration $T_n$,
anodic amplitude $A_p$ (with cathodic amplitude $A_n$ pinned by
$A_p T_p = A_n T_n$ to keep total injected charge zero), and pulse-train
frequency $f$. We record spike count, latency, and peak voltage per
configuration rather than searching for a threshold, so each grid point
is one forward simulation.

The grid has $T_p,\,T_n \in \{50,\,100,\,200,\,500\}\,\mu$s
(four values each), $A_p \in \{10,\,20,\,30,\,40,\,80\}\,\mu$A, and
pulse-train frequencies $f \in \{100,\,200,\,500,\,1000\}\,$Hz, for
$315$ valid charge-balanced forward simulations; the five invalid
combinations have $T_p + T_n$ equal to the $1000\,$Hz period. The full grid
runs in $16\,$min on the v5litepod-4 pod-slice via
\texttt{vmap}-over-configuration with hierarchical scan-checkpointing
(\cref{sec:res-throughput}). Spike crossings concentrate along the
asymmetric $T_p = 500\,\mu$s edge at low-to-mid amplitudes
($A_p \le 30\,\mu$A), and at $A_p = 80\,\mu$A for $T_p = 200\,\mu$s
(\cref{fig:biphasic-grid}); the symmetric diagonal $T_p = T_n$ is far
less effective at the same charge per pulse. The
implication for waveform design is direct: asymmetry is not a marginal
parameter --- on this cell, at this distance, swapping a symmetric
biphasic pulse for an asymmetric one with the same charge can be the
difference between sub-threshold and reliable spiking.

\begin{figure}[t]
  \centering
  \includegraphics[width=\linewidth]{figures/biphasic_grid.png}
  \caption{Charge-balanced biphasic grid on the BBP L2/3 Pyr cell at
  $d = 100\,\mu$m. Each panel is a $(T_p, T_n)$ heatmap of the maximum
  soma voltage over the four pulse-train frequencies, with one panel
  per anodic amplitude $A_p$. The dotted diagonal $T_p = T_n$ marks the
  symmetric-biphasic line; cells where $\mathit{vmax}$ crosses
  $0\,$mV correspond to a soma spike at some frequency in the sweep.
  White cells indicate configurations for which all simulated
  frequencies returned non-finite voltage traces and therefore no
  interpretable $\mathit{vmax}$; these occur mainly in the
  long-$T_p$ high-amplitude regime. Five period-filling configurations
  ($T_p = T_n = 500\,\mu$s at $1000\,$Hz) are absent from the
  forward-simulation set.}
  \label{fig:biphasic-grid}
\end{figure}

\subsection{Throughput scaling on TPU}\label{sec:res-throughput}

The point of the simulator extension is to support design-space
sweeps that were previously infeasible. We measure how it scales
across two regimes that probe different bottlenecks. A
\emph{cheap-cells} sweep isolates per-cell compile and kernel-dispatch
overhead at small problem size: \texttt{ncomp}$=2$ on the BBP Pyr
($\sim\!700$ compartments per cell after branchpoint expansion), $B$
up to $1024$ on the v5litepod-4 pod-slice. A \emph{full-fidelity}
sweep is the actual scientific workload: $50$-compartment branches on
the same Pyr cell ($\approx\!17\,500$ compartments per cell), $B$ up
to $4096$ on the same pod-slice. Each $B$ value is a single
\texttt{jax.jit(jax.vmap(\ldots))} dispatch; we report the cached
(JIT-warm) timing.

The cheap-cells curve scales linearly to $B=1024$ at $5.1\,$ms per cell
($\sim\!195$ cells/s); the full-fidelity curve scales linearly through
$B=1024$, peaks at $7.28$ cells per second at $B=2048$, and rolls over
at $B=4096$ as the four chips approach arithmetic-intensity saturation
(\cref{fig:throughput-scaling}). The full-fidelity regime depends on
the sparse-$G$ and \texttt{checkpoint\_lengths} choices in the Method
(\cref{sec:method}); without either, the $(T,\,B,\,N)$ voltage
history overflows the v5e's $16\,$GB chip memory at $B \ge 256$.

\begin{figure}[t]
  \centering
  \includegraphics[width=\linewidth]{figures/throughput_v5p4_scaling.png}
  \caption{Jaxley batch throughput on the v5litepod-4 pod-slice for two
  workload regimes: cheap cells ($700$ compartments per cell, blue) and
  full-fidelity BBP Pyr ($\sim\!17\,500$ compartments per cell, red).
  Throughput scales linearly through $B=1024$ in both regimes; the
  full-fidelity curve peaks at $7.28$ cells per second at $B=2048$ and
  saturates at $B=4096$. JIT-warm timings.}
  \label{fig:throughput-scaling}
\end{figure}

\paragraph{Comparison with NEURON.} The same $700$-compartment Pyr cell
on a $T_\text{sim} = 200\,$ms protocol runs at $0.286$ simulations per
second under serial single-core NEURON (the default BBP setup), against
$9.67$ simulations per second for Jaxley at $B = 1000$ on a single TPU
once the JIT cache is warm --- a $\sim\!34\times$ wall-clock speedup
(\cref{fig:throughput}).
NEURON has parallel-CPU variants such as
CoreNEURON~\cite{Kumbhar2019CoreNEURONA} and
experimental GPU support that we did not benchmark; the relevant
comparison for our purposes is between a CPU-serial reference and an
accelerator-batched formulation. AI accelerators are SIMT/SIMD
throughput devices, and a batched-\texttt{vmap} formulation is the
workload pattern they are designed for; a CPU-bound serial framework
cannot expose the same batch parallelism without a similar
reformulation. The practical consequence: the parameter sweeps in
\cref{sec:res-sd,sec:res-grid} finish in tens of minutes, putting
sweeps over the full waveform-and-geometry design space within reach
of a single TPU job.

\begin{figure}[t]
  \centering
  \includegraphics[width=\linewidth]{figures/bbp_throughput.png}
  \caption{Wall-clock comparison on the same $700$-compartment Pyr cell
  and $T_\text{sim} = 200\,$ms protocol: serial single-core NEURON
  (blue) versus Jaxley on a single TPU chip (red). At $B=1000$, NEURON
  takes $58.3\,$min (measured), Jaxley takes $1.7\,$min ($103.4\,$ms
  per simulation, JIT-warm), giving a $\sim\!34\times$ wall-clock
  speedup. The dashed blue line is the linear extrapolation of the
  per-NEURON-sim cost.}
  \label{fig:throughput}
\end{figure}

\subsection{Gradient-based design signal}\label{sec:res-grad}

For the differentiability we built into the pipeline to matter for
stimulation design, the gradients have to point in the right
direction --- and their direction has to mean something physical.
Establishing that gradients exist (\cref{sec:results-verify}) is a
prerequisite for, but not equivalent to, showing that they carry usable
design signal. The position gradient reported in
\cref{sec:results-verify} admits a direct physical reading. With
the electrode at $(0,\,100,\,0)\,\mu$m, the negative $y$-component of
$\partial v_\mathrm{soma,\,peak} / \partial \mathbf{x}_e$ means moving
the electrode farther from the cell along its dorsal axis reduces the
soma response, consistent with the $1/d$ decay in the field model; the
horizontal components measure the local field anisotropy induced by the
asymmetric placement of dendrites relative to the electrode. Free-form
optimisation of $w(t)$ against a spike-elicitation target under
charge-balance and amplitude constraints is a natural next step and is
left as future work (\cref{sec:limitations}). What this section
establishes is that the same pipeline that runs the empirical sweeps
of \cref{sec:res-sd,sec:res-grid} also returns design-relevant
gradients in a single backward pass --- placing direct waveform
optimisation, rather than gradient-free search, within reach of the
same compute budget.

\begin{figure}[t]
  \centering
  \includegraphics[width=0.78\linewidth]{figures/gradient_arrow.png}
  \caption{Position-gradient direction on the BBP L2/3 Pyr cell. Grey
  points are the $700$-compartment morphology projected onto the
  $y$-$z$ plane; the soma is marked with a black dot, the electrode
  ($\mathbf{x}_e = (0,\,100,\,0)\,\mu$m) with a black cross. The
  orange arrow shows the direction of
  $\partial v_\mathrm{soma,\,peak} / \partial \mathbf{x}_e$, i.e.\ the
  direction in which moving the electrode increases the soma peak.}
  \label{fig:gradient-arrow}
\end{figure}

% =============================================================
\section{Discussion}\label{sec:discussion}
% =============================================================

Genetic search over
waveforms~\cite{grill2015model,wongsarnpigoon2010energy} returns no
gradient direction for a candidate that is close to but short of an
activation target; surrogate networks trained on NEURON
outputs~\cite{hussain2024efficient} introduce an approximation error
that must be re-incurred whenever the underlying biophysical model
changes. Because the pipeline is JAX-traceable end-to-end, gradients
$\partial \ell / \partial w(t)$ of any differentiable scalar objective
$\ell$ with respect to the input waveform are returned in a single
backward pass and avoid both costs. As a
sharper receipt than the position gradient of \cref{sec:res-grad}, an
Adam loop on $\mathrm{MSE}(\hat{v}, v_\mathrm{target})$ on the full
BBP Pyr morphology recovers \texttt{gNaTs2T} from a $-15\,\%$
perturbation to within $+0.13\,\%$ in nine steps
(\texttt{fit\_demo}); the point of this fit is end-to-end gradient
flow on a realistic cell, not optimiser performance. A selectivity
Pareto sweep against a non-target cell, or a charge-minimising waveform
under a spike-elicitation constraint, follows the same template (a
differentiable spike-elicitation surrogate plus charge-balance and
amplitude constraints).

% =============================================================
\section{Limitations and future work}\label{sec:limitations}
% =============================================================

\paragraph{Limitations.}
\begin{itemize}
  \item \emph{Volume conductor.} The field model is a homogeneous,
  isotropic, infinite point-source monopole. Tissue anisotropy,
  finite electrode geometry (cuff, paddle, contact size), and the
  bone--CSF--grey-matter layering used in DBS planning are not
  represented.
  \item \emph{Quasi-static assumption.} Tissue capacitance is ignored;
  this is standard at neural-stimulation frequencies but breaks down at
  very short pulse edges.
  \item \emph{Float-32 caveat.} The production sweeps run in float-32
  after a $0.77\,\%$ threshold-difference check against float-64 on one
  HH cable configuration; the check should be repeated per cell type
  before reusing float-32 in a regime we have not validated.
  \item \emph{PV cell stability.} The full-morphology BBP L2/3 LBC
  basket cell is unstable in Jaxley's default integrator (NaN at
  $\sim\!30\,$ms even unstimulated), traced to BBP's
  $g_\mathrm{pas} = 1\!\times\!10^{-6}\,$S/cm$^2$ on PV apical and
  basal --- $30\times$ lower than the Pyr value and tolerated by NEURON
  but not by the Jaxley defaults. Single-compartment PV parity passes
  (\cref{sec:res-neuron}); all full-morphology runtime claims in this
  paper are therefore Pyr-only.
  \item \emph{Pyr Ih path-distance traversal.} The
  $-0.38\,$mV resting-potential offset against NEURON
  (\cref{sec:res-neuron}) is the residual of approximating BBP's
  path-distance apical $I_h$ gradient with Euclidean distance.
  Replacing the approximation with a parent-pointer breadth-first
  traversal would close the residual but we did not implement it.
  \item \emph{Solver floor.} \texttt{checkpoint\_lengths} bounds memory
  but not arithmetic. The backward-Euler scan dominates the per-step
  cost on a v5e chip and is what saturates the throughput curve at
  $B=2048$ in \cref{sec:res-throughput}.
\end{itemize}

\paragraph{Future work.}
Three lines of follow-up sit immediately on top of what is here. First,
restoring full PV cell coverage by stabilising the LBC integration ---
either with an explicit smaller-timestep regime, a leak-conductance
floor, or a stiff-aware solver --- and applying the BBP path-distance
$I_h$ gradient to PV apical and basal as in Pyr. Second, gradient-based
optimisation of the free-form waveform $w(t)$ against a
spike-elicitation surrogate, under charge-balance and amplitude
constraints, evaluated against the strength--duration baseline of
\cref{sec:res-sd}. Third, selectivity: two cell models with different
geometries placed in the same extracellular field, optimising $w(t)$
to drive a soma spike in the target cell while keeping the non-target
subthreshold. The Pareto frontier between activation and sparing is
exactly what the differentiable pipeline is built to trace; a
penalty-based formulation
$\mathcal{L}(w) = \lVert w \rVert_1 \Delta t
+ \lambda\,\tilde{F}(w;\theta_2)$ subject to
$\tilde{F}(w;\theta_1) \ge \tau$ keeps the problem smooth and lets a
$\lambda$-sweep trace the frontier directly.


% =============================================================
\section{Conclusion}\label{sec:conclusion}
% =============================================================

We have added extracellular stimulation to Jaxley by recovering the
discrete activating function from the simulator's own axial diffusion
operator and injecting it as an equivalent compartmental current through
the public stimulus API. The pipeline composes with \texttt{jit},
\texttt{vmap}, \texttt{grad}, and \texttt{jax.sharding} without a fork
of the simulator, agrees with Jaxley's internal voltage vector field to
$\sim\!10^{-13}$ relative error, and runs the BBP L2/3 pyramidal cell at
full $50$-compartment-per-branch fidelity on a four-chip v5e pod-slice.
What this enables, and what we leave to follow-up work, is gradient-based
search over the waveform itself --- the design problem that motivated
the extension.



\bibliographystyle{plain}
\bibliography{ref}

\appendix

% =============================================================
\section{Branchpoint elimination}\label{app:branchpoint-elimination}
% =============================================================

Jaxley's internal representation introduces branchpoint pseudo-nodes
at cable junctions to enforce Kirchhoff's current law. Partitioning
$G$ into real ($R$) and branchpoint ($B$) blocks gives
\[
  \begin{pmatrix} G_{RR} & G_{RB} \\ G_{BR} & G_{BB} \end{pmatrix}
  \begin{pmatrix} \mathbf{v}_R \\ \mathbf{v}_B \end{pmatrix}
  =
  \begin{pmatrix} [G\mathbf{v}]_R \\ [G\mathbf{v}]_B \end{pmatrix}.
\]
Branchpoint nodes carry no membrane, so $c_b A_b = 0$ at $b$ and the
cable equation at $b$ reduces to the algebraic constraint
$[G\mathbf{v}]_B = 0$, giving $\mathbf{v}_B = -G_{BB}^{-1}G_{BR}\mathbf{v}_R$
and the Schur complement
\[
  \tilde{G} \;=\; G_{RR} - G_{RB}\,G_{BB}^{-1}\,G_{BR}.
\]
The same operator applies to the extracellular forcing:
$\tilde{G}\boldsymbol{\phi}_R$ gives the activating function at the
real compartments, where $\boldsymbol{\phi}_R$ is the field sampled at
real compartment centres only. The reduced operator acts only on real
compartments, so no $\boldsymbol{\phi}_B$ value enters the forcing. Our implementation delegates this assembly to a
private Jaxley helper rather than reconstructing it from scratch, so
the operator stays in sync with upstream solver changes.


% =============================================================
\section{Equivalent-current encoding}\label{app:equivalent-current}
% =============================================================

The activating function $G\boldsymbol{\phi}$ has units mV/ms (a
voltage rate), but Jaxley's stimulus API accepts current in nA.
Internally Jaxley converts
$i_{\text{ext}}\,[\mu\text{A/cm}^2] = i_{\text{nA}}\,[\text{nA}] /
A\,[\mu\text{m}^2]\cdot 10^5$ and adds $i_{\text{ext}}/c$ to $dv/dt$.
To inject $[G\boldsymbol{\phi}]_j$ on the other side of this
conversion, we encode
\[
  i_{\text{ecs},j}\,[\text{nA}] \;=\;
    c_j \cdot [G\boldsymbol{\phi}(t)]_j \cdot \frac{A_j}{10^5}.
\]
The $c_j$ and $A_j$ factors cancel in the round trip; they are an
encoding artifact of using the public API rather than modifying the
solver. In matrix form,
\[
  \mathbf{I}_{\text{ecs}} \;=\;
    \left(\mathbf{c} \odot \frac{\mathbf{A}}{10^5}\right)
    \odot (G\boldsymbol{\Phi})
    \;\in\; \mathbb{R}^{N\times T}.
\]


% =============================================================
\section{BBP cell-model implementation notes}\label{app:bbp-implementation}
% =============================================================

Two implementation issues affect reproducibility of the BBP cell models
in Jaxley but are unrelated to the paper's results.

\paragraph{SWC soma conversion.} The BBP source morphologies represent
the soma as a two-dimensional contour, which the SWC format cannot
encode. The Neurolucida-to-SWC conversion tool emits soma points with
$r = 0$, which Jaxley turns into degenerate zero-area compartments
that produce NaN voltages at the first integration step. We patch the
converted SWC files at load time so the soma points carry a
non-degenerate radius.

\paragraph{Calcium-channel insertion order.} Calcium signalling is
sensitive to channel insertion order. The calcium Nernst reversal
must be established before any calcium-carrying channel, which in
turn must precede the calcium pump and the calcium-activated
potassium channel.
\texttt{cell\_factory.\_insert\_calcium\_chain} enforces
\texttt{CaNernstReversal} $\to$ \texttt{CaHVA/CaLVA} $\to$
\texttt{CaPump} $\to$ \texttt{SKE2}; out-of-order insertion produces
uninitialised state references at integration time.


% =============================================================
\section{Apical \texorpdfstring{$I_h$}{Ih} distance-dependent
gradient}\label{app:ih-gradient}
% =============================================================

The \texttt{cADpyr229} apical $I_h$ conductance follows BBP's
distance-dependent gradient
\[
  \bar{g}_{I_h}(d)
    = \bigl(-0.8696 + 2.087\,e^{0.0031\,d}\bigr)
      \cdot 8\!\times\!10^{-5}\;\mathrm{S/cm}^2,
\]
where $d$ is path distance from the soma centre in $\mu$m. We
approximate $d$ with the Euclidean distance to the soma centre, which
is close for roughly-vertical apical trees. Before the fix the
resting-potential offset against NEURON was $-0.62\,$mV; the
Euclidean approximation closes it to $-0.38\,$mV, and a path-distance
traversal would close it further (\cref{sec:limitations}).


% =============================================================
\section{Reproducibility}\label{app:reproducibility}
% =============================================================

The build pipeline binds the entire compute environment, not just the
Python packages. A \texttt{requirements.txt} inside a Dockerfile pins
Python packages and a base OS image but leaves the C and
C\nolinkurl{++} compilers, the CUDA toolchain, the Python interpreter
build, and the cloud topology unspecified --- each of which can drift
between machines and change numerical results. The stack we use
closes those gaps.

\paragraph{Python dependencies (\texttt{uv}).} A Rust-implemented
package manager that resolves direct and transitive Python
dependencies and records every resulting wheel and its hash in a
lockfile (\texttt{uv.lock}). Two machines sharing the lockfile
install the same versions of every Python package, regardless of
install order.

\paragraph{System environment (Nix).} A Nix flake (\texttt{flake.nix})
specifies the Python interpreter (3.11), CUDA runtime, JAX, the C
compiler that built them, and command-line tooling. Store paths are
determined by the hash of the full transitive build graph (pinned via
\texttt{flake.lock}), so the same flake on the same target system
produces the same dependency closure.

\paragraph{Cloud infrastructure (OpenTofu).} The TPU pod-slice, Cloud
SQL Postgres database, and GCS artifact bucket are declared as
OpenTofu configuration in \texttt{infra/tofu/}; \texttt{tofu apply}
brings the topology up, \texttt{tofu destroy} tears it down.

\paragraph{Reproducing the figures.} A reader entering the Nix shell
gets a deterministic Python and CUDA environment; the project
\texttt{README} carries the exact commands. Reproducing the figures
from pinned data does not require any TPU-scale workload.

\paragraph{Artifact set.} The curated data backing every figure ---
parity \texttt{.npz} files, the strength--duration and biphasic-grid
Zarr stores, the throughput sweeps, and the gradient receipt ---
totals a few megabytes and is published as a separate artifact
distribution rather than committed to the source repository, available
at \url{https://github.com/Marco-Christiani/jaxley-extracellular/releases/tag/paper-v0.1}.


\end{document}
